00735: \section*{Pr\'eface}
00736: \addcontentsline{toc}{section}{Pr\'eface}
00737: 
00738: La pr\'esente \'edition du S\'eminaire de G\'eom\'etrie Alg\'ebrique 6
00739: reproduit \`a peu de choses pr\`es la premi\`ere \'edition, distribu\'ee par
00740: l'I.H.E.S. sous forme d'expos\'es multigraphi\'es. Les corrections effec-
00741: tu\'ees portent essentiellement sur des fautes d'impression, sauf dans
00742: les expos\'es I et III, o\`u de l\'eg\`eres modifications ont \'et\'e apport\'ees
00743: par L. Illusie \`a quelques \'enonc\'es incorrects dans la version primiti-
00744: ve. D'autre part, des notes ont \'et\'e rajout\'ees en bas de page, en par-
00745: ticulier dans l'expos\'e XIV par A. Grothendieck, afin de signaler des
00746: progr\`es r\'ecents dans des questions encore ouvertes lors du S\'eminaire.
00747: Signalons enfin que l'expos\'e XI, de D. Ferrand, concernant la th\'eorie
00748: du d\'eterminant pour les complexes parfaits, n'a pas \'et\'e r\'edig\'e, et ne
00749: sera donc pas publi\'e; pour \'eviter des erreurs de r\'ef\'erence, nous
00750: avons gard\'e telle quelle la num\'erotation des expos\'es suivants.
00751: 
00752: \begin{flushright}
00753: Juin 1971
00754: \end{flushright}
00755: 
00756: 
00757: \section*{Introduction}
00758: \addcontentsline{toc}{section}{Introduction}
00759: 
00760: Le but du pr\'esent S\'eminaire est de d\'evelopper une th\'eorie globale des
00761: intersections, et une formule de Riemann-Roch, pour des sch\'emas quelconques.
00762: Le lecteur trouvera une description du programme du S\'eminaire dans l'expos\'e 0.
00763: La possibilit\'e en principe d'une d\'emonstration d'une formule de Riemann-Roch
00764: pour les sch\'emas r\'eguliers g\'en\'eraux, suivant les lignes du rapport bien con-
00765: nu de Borel-Serre en 1958, \'etait claire d\`es ce moment, du moins pour un mor-
00766: phisme projectif. L'extension au cas g\'en\'eral est moins \'evidente; le program-
00767: me dans lequel elle s'ins\`ere (et partiellement r\'ealis\'e dans notre s\'eminaire)
00768: remonte \`a 1960. Comme c'\'etait \'egalement le cas pour la th\'eorie de dualit\'e
00769: des faisceaux coh\'erents (cf. R. Hartshorne, \textit{Residues and Duality}, Lecture
00770: Notes in Mathematics no 20, Springer), un outil essentiel pour formuler une
00771: t\'eorie satisfaisante est la th\'eorie des cat\'egories d\'eriv\'ees de Verdier,
00772: dont la connaissance est indispensable pour la compr\'ehension du S\'eminaire.
00773: 
00774: A part cette th\'eorie, l'\'etude du S\'eminaire n'exige gu\`ere qu'une connaissance
00775: g\'en\'erale des fondements de la G\'eom\'etrie Alg\'ebrique, tels qu'ils sont expos\'es
00776: dans EGA, chapitres I, II, III; nous n'aurons en plus que des r\'ef\'erences oc-
00777: casionnelles \`a faire \`a EGA IV, pour certains faits concernant les morphismes
00778: plats ou lisses, qui pour l'essentiel figurent \'egalement dans les premiers
00779: expos\'es de SGA 1. Enfin, pour d\'evelopper certains r\'esultats avec toute la
00780: g\'en\'eralit\'e souhaitable pour les applications, nous faisons usage parfois de
00781: la notion de site annel\'e et de topos annel\'e, pour laquelle nous renvoyons \`a
00782: SGA 4, expos\'es I \`a IV. Le lecteur qui ignorerait le langage des sites et topos
00783: pourra remplacer partout lesdits objets par des espaces topologiques ordinai-
00784: res, les objets du topos \'etant alors remplac\'es par des ouverts de ces espaces;
00785: mais nous lui conseillons n\'eanmoins, de pr\'ef\'erence, de s'assimiler le langage
00786: des topos, qui fournit un principe d'unification extr\^emement commode.
00787: 
00788: La notion fondamentale pour la th\'eorie pr\'esent\'ee ici est celle de
00789: complexe de Modules parfait, et ses diverses variantes, d\'evelopp\'ees dans les
00790: expos\'es I, II, III. Il semble clair que ces notions, importantes \'egalement
00791: dans d'autres contextes en G\'eom\'etrie Alg\'ebrique (notamment pour les formules
00792: de type Lefschetz-Verdier en cohomologie \`a coefficients discrets (SGA 5)
00793: ou coh\'erents), auront aussi leur r\^ole \`a jouer en dehors de la G\'eom\'etrie
00794: Alg\'ebrique, notamment pour la formulation d'une variante analytique com-
00795: plexe du th\'eor\`eme de Riemann-Roch pr\'esent\'e ici, ou de variantes convenables
00796: du th\'eor\`eme de l'index d'Atiyah-Singer. Quelques indications dans ce sens
00797: seront donn\'ees dans Exp. II. Malheureusement, il manque encore \`a l'heure
00798: actuelle un \'enonc\'e, m\^eme heuristique, qui engloberait ces deux th\'eor\`emes
00799: (dont le premier pour l'instant reste conjectural). Il manque apparemment
00800: une id\'ee nouvelle, comme il en manque aussi en G\'eom\'etrie Alg\'ebrique pour
00801: parvenir \`a une d\'emonstration du th\'eor\`eme de Riemann-Roch en dehors d'hypo-
00802: th\`eses projectives (cf. Exp. XIV, no 2).
00803: 
00804: Nous n'aurons \`a faire nul usage dans ce s\'eminaire de la th\'eorie locale
00805: des intersections, expos\'ee dans J.P. Serre, \textit{Algebre Locale. Multiplicites}
00806: (Lectures Notes in Mathematics no 11, Springer), et nous nous contenterons sim-
00807: plement de signaler ici que ce cours de Serre a eu une influence \'evidente
00808: sur la gen\`ese de la th\'eorie en 1957. Nous n'utiliserons pas non plus la
00809: th\'eorie de l'anneau de Chow d\'evelopp\'ee dans Seminaire Chevalley 1958, expo-
00810: s\'es 2 et 3 (Ecole Normale Superieure). Cette th\'eorie est li\'ee de facon es-
00811: sentielle \`a des hypoth\`eses de non singularit\'e, alors que le but de notre
00812: s\'eminaire est au contraire de d\'evelopper une th\'eorie des intersections
00813: sur les sch\'emas g\'en\'eraux (et m\^eme les topos localement annel\'es g\'en\'eraux)
00814: voyez \`a ce sujet les commentaires dans Exp. XIV no 4 et no 8, donnant les
00815: relations entre notre th\'eorie et celle de l'anneau de Chow, et posant la
00816: question d'une g\'en\'eralisation de cette derni\`ere. Nous pouvons dire que
00817: le Seminaire pr\'esente une th\'eorie des intersections coherente et se suffi-
00818: sant \`a elle m\^eme, mais nullement exhaustive des diff\'erents points de vue
00819: utiles (voire indispensables) en th\'eorie des intersections, et qu'il con-
00820: vient par suite de le completer par les expos\'es cites de Serre et de
00821: Chevalley.
00822: 
00823: Nous avons joint au s\'eminaire (en Appendice \`a Exp. 0) le rapport
00824: Grothendieck de 1957 ``Classes de Faisceaux et Theoreme de Riemann-Roch'',
00825: qui avait servi de base au s\'eminaire de Borel et Serre \`a Princeton la
00826: m\^eme annee, ainsi qu'\`a leur article deja cite. Ce rapport esquisse deux
00827: demonstrations du theoreme de Riemann-Roch pour les varietes quasi-projec-
00828: tives non singuli\`eres, dont la premi\`ere, valable pour le moment seulement
00829: en caracteristique nulle, mais donnant en revanche un resultat un peu plus
00830: precis dans le cas d'une immersion, ne figure pas encore dans la littera-
00831: ture (sans exclure le travail de Borel-Serre ni le present s\'eminaire). La
00832: lecture de ce rapport ne presuppose bien entendu celle d'aucun autre expose
00833: du s\'eminaire, et peut meme servir d'introduction \`a l'\'etude de ce dernier
00834: au meme titre que l'expose 0, surtout pour un lecteur qui ne serait pas
00835: encore familier avec Borel-Serre. La demonstration \`a laquelle nous venons
00836: de faire allusion s'applique d'ailleurs essentiellement telle quelle au cas
00837: d'un morphisme projectif d'espaces analytiques complexes, et s'appliquera
00838: sans doute egalement en caracteristique quelconque, une fois resolu le pro-
00839: bleme des operations ``puissances exterieures'' dans la categorie derivee
00840: (cf. Exp. XIV, nos 1,2,3). C'est l'absence d'une etude de cette question
00841: qui constitue sans doute la lacune la plus serieuse dans ce Seminaire, dans
00842: l'optique meme ou nous nous y sommes places.
00843: 
00844: On remarquera l'absence, dans la table des matieres du present Semi-
00845: aire, de deux des participants mentionnes sur la page de garde. L'un, J.P.
00846: Jouanolou, avait pris une part active a l'elaboration technique de la pre-
00847: miere partie du seminaire, mais avait ete empeche de prendre part aux expo-
00848: ses oraux; on lui doit notamment l'assertion d'independance lineaire conte-
00849: nue dans l'important enonce VI 1.1. donnant la structure de l'anneau $K^{\bullet}$
00850: des classes de Modules localement libres sur un fibre projectif (qui n'etait prouve
00851: auparavant que lorsqu'on supposait que le schema de base admet un Module
00852: inversible ample). L'autre, J.P. Serre, avait fait deux exposes oraux,
00853: logiquement independants du reste du Seminaire, et qu'il a par suite prefere
00854: publier sous forme d'article separe (cf. J.P. Serre, \textit{Groupes de
00855: Grothendieck des schemas en groupes reductifs deployes} a paraitre dans
00856: \textit{Publications Mathematiques}, no 34).
00857: 
00858: Comme dans les autres parties du Seminaire de Geometrie Algebrique du
00859: Bois Marie, les sigles SGA 1 a SGA 6 renvoient aux differentes parties
00860: dudit Seminaire, le chiffre romain suivant ce sigle indiquant le no de l'ex-
00861: pose, et les chiffres arabes qui le suivent correspondant a la numerotation
00862: interieure de l'expose en question; pour les references interieures au pre-
00863: sent Seminaire, on omet le sigle SGA 6, en utilisant des references style
00864: I 4.9 (signifiant: expose I, enonce 4.9). Le sigle EGA refere aux \textit{Elements
00865: de Geometrie Algebrique} de J. Dieudonne et A. Grothendieck.
00866: 
00867: 
00868: \section*{Table des matieres}
00869: \addcontentsline{toc}{section}{Table des matieres}
00870: 
00871: \begin{center}
00872: \textbf{EXPOSE 0}
00873: 
00874: \medskip
00875: \textit{Esquisse d'un Programme pour une Theorie des Intersections\\
00876: sur les Schemas Generaux}\\
00877: par A. Grothendieck \dotfill 1
00878: 
00879: \bigskip
00880: \textbf{APPENDICE}
00881: 
00882: \medskip
00883: \textit{Classes de Faisceaux et Theoreme de Riemann-Roch}\\
00884: par A. Grothendieck \dotfill 20
00885: 
00886: \bigskip
00887: \textbf{EXPOSE I}
00888: 
00889: \medskip
00890: \textit{Generalites sur les Conditions de Finitude dans les Categories\\
00891: Derivees}\\
00892: par L. Illusie \dotfill 56
00893: 
00894: \bigskip
00895: \textbf{EXPOSE II}
00896: 
00897: \medskip
00898: \textit{Existence de Resolutions Globales}\\
00899: par L. Illusie \dotfill 160
00900: 
00901: \bigskip
00902: \textbf{EXPOSE III}
00903: 
00904: \medskip
00905: \textit{Conditions de Finitude Relatives}\\
00906: par L. Illusie \dotfill 190
00907: 
00908: \bigskip
00909: \textbf{EXPOSE IV}
00910: 
00911: \medskip
00912: \textit{Groupes de Grothendieck des Topos Anneles}\\
00913: par L. Illusie \dotfill 222
00914: 
00915: \bigskip
00916: \textbf{EXPOSE V}
00917: 
00918: \medskip
00919: \textit{Generalites sur les $\lambda$-Anneaux}\\
00920: par P. Berthelot \dotfill 297
00921: 
00922: \bigskip
00923: \textbf{EXPOSE VI}
00924: 
00925: \medskip
00926: \textit{Le $K^{\bullet}$ d'un Fibre Projectif: Calculs et Consequences}\\
00927: par P. Berthelot \dotfill 365
00928: \end{center}
00929: 
00930: 
00931: \begin{center}
00932: \textbf{EXPOSE VII}
00933: 
00934: \medskip
00935: \textit{Immersions Regulieres et Calcul du $K^{\bullet}$ d'un Schema Eclate}\\
00936: par P. Berthelot \dotfill 416
00937: 
00938: \bigskip
00939: \textbf{EXPOSE VIII}
00940: 
00941: \medskip
00942: \textit{Le Theoreme de Riemann-Roch}\\
00943: par P. Berthelot \dotfill 466
00944: 
00945: \bigskip
00946: \textbf{EXPOSE IX}
00947: 
00948: \medskip
00949: \textit{Quelques Calculs de Groupes $K_{\bullet}$}\\
00950: par P. Berthelot \dotfill 498
00951: 
00952: \bigskip
00953: \textbf{EXPOSE X}
00954: 
00955: \medskip
00956: \textit{Formalisme des Intersections sur les Schemas Algebriques Propres}\\
00957: par O. Jussila\\
00958: Avec un Appendice par A. Grothendieck
00959: 
00960: \medskip
00961: \textit{Specialisation en Theorie des Intersections} \dotfill 595
00962: 
00963: \bigskip
00964: \textbf{EXPOSE XI -- Non redige}
00965: 
00966: \bigskip
00967: \textbf{EXPOSE XII}
00968: 
00969: \medskip
00970: \textit{Un Theoreme de Representabilite Relative sur le Foncteur de Picard}\\
00971: par M. Raynaud (redige par S. Kleiman) \dotfill 595
00972: 
00973: \bigskip
00974: \textbf{EXPOSE XIII}
00975: 
00976: \medskip
00977: \textit{Les Theoremes de Finitude pour le Foncteur de Picard}\\
00978: par S. Kleiman \dotfill 616
00979: 
00980: \bigskip
00981: \textbf{EXPOSE XIV}
00982: 
00983: \medskip
00984: \textit{Problemes Ouverts en Theorie des Intersections}\\
00985: par A. Grothendieck \dotfill 667
00986: 
00987: \bigskip
00988: \textit{Index Terminologique} \dotfill 707
00989: 
00990: \medskip
00991: \textit{Index des Notations} \dotfill 710
00992: \end{center}
00993: 
00994: 
00995: \section*{Expose 0}
00996: \addcontentsline{toc}{section}{Expose 0 -- Esquisse d'un Programme}
00997: 
00998: \begin{center}
00999: \Large\textbf{ESQUISSE D'UN PROGRAMME POUR UNE THEORIE DES INTERSECTIONS}\\[0.3cm]
01000: \Large\textbf{SUR LES SCHEMAS GENERAUX}\\[0.5cm]
01001: \textit{par A. Grothendieck}
01002: \end{center}
01003: 
01004: \bigskip
01005: 
01006: Le present expose est de nature introductif, et sa lecture n'est pas
01007: logiquement indispensable pour l'etude du Seminaire. Il s'adresse plus par-
01008: ticulierement aux lecteurs au courant de la version provisoire du theoreme
01009: de Riemann-Roch, telle qu'elle est exposee dans le rapport de Borel-Serre
01010: [2], ou dans le rapport de Grothendieck deja mentionne dans l'Introduction
01011: (cite [RRR], par la suite), et qui est reproduit sous forme d'appendice a
01012: la fin du present expose.
01013: 
01014: \subsection*{1. Rappelons la formule de Riemann-Roch}
01015: \addcontentsline{toc}{subsection}{1. Rappel}
01016: 
01017: Rappelons la formule de Riemann-Roch pour un morphisme propre
01018: $f: X \to Y$
01019: de schemas quasi-projectifs et lisses sur un corps $k$, et un faisceau co-
01020: herent $F$ sur $X$, dont la classe dans le groupe $K(X)$ des classes de fais-
01021: ceaux coherents sur $X$ est notee par $\Cl(F)$:
01022: \begin{equation}\label{eq:RR1}
01023: \Todd(T_Y) \cdot \ch_Y(f_{!}(\Cl(F))) = f_{*}(\Todd(T_X) \cdot \ch_X(F))
01024: \end{equation}
01025: Cette formule est valable dans $A(Y) \otimes \mathbf{Q}$, ou $A(Y)$ est l'anneau de Chow de $Y$,
01026: le $f_{*}$ du second membre etant deduit par tensorisation par $\mathbf{Q}$ de l'homomor-
01027: phisme ``image directe de cycles''
01028: \[
01029: f_{*}: A(X) \longrightarrow A(Y),
01030: \]
01031: celui du premier etant la caracteristique d'Euler-Poincare de $F$ relative-
01032: ment a $f$:
01033: \[
01034: f_{!}(\Cl(F)) = \sum (-1)^i \Cl(R^i f_{*}(F));
01035: \]
01036: $\ch_X$, $\ch_Y$ designent les caracteres de Chern sur $X$ resp. sur $Y$, et $T_X$, $T_Y$
01037: les fibres tangents a $X$ resp. a $Y$. Comme on sait, $\Todd(-)$ et $\ch(-)$ sont
01038: des polynomes universels a coefficients dans $\mathbf{Q}$ en les classes de Chern de
01039: l'argument. Le terme constant de $\Todd(-)$ etant $1$, c'est un element inversi-
01040: ble pour toute valeur de l'argument, de sorte que la formule (\ref{eq:RR1}) prend
01041: apres multiplication par $\Todd(T_Y)^{-1}$ la forme equivalente, plus commode pour
01042: notre propos:
01043: \begin{equation}\label{eq:RR2}
01044: \ch_Y(f_{!}(\Cl(F))) = f_{*}\bigl(\Todd(T_f) \cdot \ch_X(F)\bigr);
01045: \end{equation}
01046: ou on a pose
01047: \begin{equation}\label{eq:Tf}
01048: T_f = T_X - f^{*}(T_Y) \in K(X).
01049: \end{equation}
01050: 
01051: \subsection*{2. Examens des generalisations}
01052: 
01053: De sorte que $T_f$ joue le role d'un fibre tangent relatif virtuel de $X$ sur
01054: $Y$. Dans le cas ou le morphisme $f$ est lisse (i.e. a application tangente
01055: partout surjective), on a simplement
01056: \[
01057: T_f = T_{X/Y} \quad \text{(fibre tangent le long des fibres),}
01058: \]
01059: tandis que dans le cas ou $f: X \to Y$ est une immersion, on trouve
01060: \[
01061: T_f = -N_{X/Y}
01062: \]
01063: ou $N_{X/Y}$ designe le faisceau normal de $X$ dans $Y$.
01064: 
01065: Un des buts principaux du Seminaire est de generaliser (\ref{eq:RR2}) simulta-
01066: nement dans deux directions:
01067: 
01068: \smallskip
01069: \textbf{a)} Se debarrasser de l'hypothese de l'existence d'un corps de base $k$.
01070: 
01071: \smallskip
01072: \textbf{b)} Remplacer les hypotheses de regularite sur $Y, X$ par une hypothese
01073: de ``regularite locale'' de $f$.
01074: 
01075: \smallskip
01076: Enfin, chemin faisant nous nous occuperons egalement du probleme:
01077: 
01078: \smallskip
01079: \textbf{c)} Eliminer les hypotheses de quasi-projectivite qui, en l'absence
01080: d'un corps de base, s'expriment par l'existence de Modules inversibles am-
01081: ples sur $X$ et sur $Y$.
01082: 
01083: \subsubsection*{2.1. Generalisation a)}
01084: 
01085: Examinons d'abord la generalisation a), en gardant cependant les hypo-
01086: theses de regularite et d'existence de Modules inversibles amples sur $X$
01087: et sur $Y$.
01088: 
01089: La definition de $K(X), K(Y)$ et de l'homomorphisme $f_{!}: K(X) \to K(Y)$
01090: n'offre alors pas de nouveau probleme, grace au fait que $X$ et $Y$ sont regu-
01091: liers. La voie la plus naturelle pour donner un sens a (\ref{eq:RR2}) semble donc
01092: consister en la definition d'anneaux de Chow $A(X), A(Y)$ et d'un homomorphis-
01093: me de groupes
01094: \[
01095: f_{*}: A(X) \longrightarrow A(Y),
01096: \]
01097: en l'etablissement d'une theorie des classes de Chern, fournissant des ap-
01098: plications
01099: \[
01100: c_i: K(X) \longrightarrow A(X),
01101: \]
01102: et de meme pour $Y$, et enfin en la description d'un element fibre tangent
01103: relatif virtuel
01104: \[
01105: T_f \in K(X).
01106: \]
01107: 
01108: \subsubsection*{2.2. Le fibre tangent relatif virtuel}
01109: 
01110: Pour ce qui est de cette derniere, une definition s'offre de facon
01111: assez evidente. On utilise le fait que le morphisme $f$, grace a l'hypothese
01112: d'existence d'un Module inversible ample sur $X$ et sur $Y$, peut se factoriser
01113: en
01114: \begin{equation}
01115: X \xrightarrow{i} X' \xrightarrow{f'} Y,
01116: \end{equation}
01117: ou $i$ est une immersion fermee, et $f'$ un morphisme lisse; par exemple on
01118: pourra prendre $X' = \mathbf{P}^r_Y$ pour $r$ convenable. On posera alors
01119: \begin{equation}
01120: T_f = i^{*}(T_{X'/Y}) - \check{N}_{X/X'},
01121: \end{equation}
01122: ou $\Omega^1_{X'/Y}$ est le Module localement libre des 1-differentielles relatives de
01123: $X'$ sur $Y$ (ou Module cotangent relatif de $X'$ sur $Y$), et le Module conor-
01124: mal de $X$ dans $X'$, egal par definition a $J/J^2$, ou $J$ est l'Ideal sur $X'$ defi-
01125: nissant le sous-schema ferme $X$; comme $X$ et $X'$ sont reguliers, on en con-
01126: clut d'ailleurs que $N_{X/X'}$ est egalement localement libre. Enfin $\check{N}$ designe
01127: le dual de $N$. On verifie sans peine (Exp. VIII) que l'element $T_f$ defini
01128: par la formule ci-dessus ne depend pas de la factorisation de $f$ choisie.
01129: 
01130: \subsubsection*{2.3. La filtration et l'anneau gradue}
01131: 
01132: Quant a une definition, sous les conditions envisagees, d'un anneau
01133: de Chow et d'une theorie des classes de Chern correspondante, bien qu'il
01134: soit plausible qu'une telle theorie doive pouvoir se developper sur le modele
01135: de celle deja connue pour les schemas algebriques lisses et quasi-projec-
01136: tifs, il semble que ce travail n'ait pas ete fait a l'heure actuelle, et a
01137: fortiori il n'a pas ete entrepris dans le Seminaire. On introduit par contre
01138: un autre anneau, qui jouera un role analogue a celui de l'anneau de Chow, et
01139: dont le produit tensoriel par $\mathbf{Q}$ est effectivement isomorphe a $A(X) \otimes_{\mathbf{Z}} \mathbf{Q}$ dans
01140: le cas ou $X$ est quasi-projectif et lisse sur un corps $k$. Cet anneau est
01141: l'anneau gradue associe a $K(X)$, pour une filtration convenable de cet anneau.
01142: 
01143: Une premiere filtration naturelle de $K(X)$, qui est celle envisagee dans
01144: [RRR], est la ``filtration topologique'' par la codimension du support des
01145: faisceaux coherents, $\operatorname{Fil}^i K(X)$ etant engendre par les $\operatorname{Cl}(F)$ avec $F$ cohe-
01146: rent et $\operatorname{codim}(\operatorname{supp} F, X) \geq i$. Une premiere question est celle de la compatibi-
01147: lite de cette filtration avec la structure d'anneau, question resolue dans
01148: [RRR], dans le cas $X$ lisse sur un corps $k$, grace au ``moving lemma'' de Chow.
01149: Ce point admis, on trouve donc un anneau gradue associe, que nous noterons
01150: ici $\Gr^*_{\mathrm{top}}(X)$. Dans [RRR] on definit un homomorphisme d'anneaux surjectif
01151: \[
01152: \phi: A(X) \longrightarrow \Gr^*_{\mathrm{top}}(X),
01153: \]
01154: (transformant la classe d'un cycle premier $Z$ en $\Cl_X(\mathcal{O}_Z)$,) dont il s'avere
01155: ensuite (Exp. XIV 4.2) que le noyau est un sous-groupe de torsion. De plus,
01156: on montre dans [RRR] comment les images dans $\Gr^*_{\mathrm{top}}(X)$ des classes de Chern
01157: $c_i(F)$ peuvent se calculer directement en termes de $\Cl(F) \in K(X)$ et des ope-
01158: rations de puissance exterieure dans $K(X)$.
01159: 
01160: Ceci montre que, pour $X$ schema noetherien regulier muni d'un Module
01161: inversible ample, l'anneau $\Gr^*_{\mathrm{top}}(X) \otimes \mathbf{Q}$ est un substitut commode pour l'anneau
01162: de Chow hypothetique $A(X)$; il a de plus l'avantage sur ce dernier que c'est
01163: directement dans cet anneau que se font certains calculs de [RRR] aboutis-
01164: sant a des formules importantes dans cet anneau, formules qui meme pour $X$
01165: projectif et lisse sur un corps ne sont pas etablies a l'heure actuelle dans
01166: $A(X)$ lui-meme (Exp. XIV, 4.3 et 4.4). Pour justifier entierement ce point
01167: de vue, il faudrait cependant resoudre le probleme signale plus haut de la
01168: compatibilite de la filtration topologique de $K(X)$ avec sa structure d'an-
01169: neau, probleme qui semble essentiellement lie au ``moving lemma'' de Chow, qui
01170: lui-meme est l'ingredient technique essentiel de la theorie de l'anneau de
01171: Chow que nous avons pretendu pouvoir eviter.
01172: 
01173: \subsubsection*{2.4. La $\lambda$-filtration}
01174: 
01175: Pour eviter ce probleme, on munit $K(X)$ d'une deuxieme filtration,
01176: plus fine que la filtration topologique, et qui peut se decrire en termes
01177: purement algebriques a partir de la structure de $\lambda$-anneau augmente de $K(X)$.
01178: La definition de cette filtration est suggeree par l'expression donnee dans
01179: [RRR] des classes de Chern d'un $x \in K(X)$ dans $\Gr^*_{\mathrm{top}}(X) \otimes \mathbf{Q}$:
01180: \begin{equation}
01181: c_i(x) = \gamma^i(x - \varepsilon(x)) \mod \Fil^{i+1}_{\mathrm{top}} K(X),
01182: \end{equation}
01183: ou $\gamma^i$ designe une certaine combinaison lineaire des $\lambda^j$ ($j \leq i$) explicitee dans
01184: loc. cit., $\varepsilon$ etant l'augmentation. On montre que $\gamma^i(x - \varepsilon(x))$ est bien de
01185: filtration topologique $\geq i$, ce qui donne un sens a la formule ci-dessus. On defi-
01186: nit alors la $\lambda$-filtration de $K(X)$ par les ideaux $\Fil^i_{\lambda} K(X)$, ou $\Fil^i_{\lambda} K(X)$
01187: est forme par les sommes d'expressions de la forme
01188: \[
01189: \gamma^{i_1}(x_1 - \varepsilon(x_1)) \cdots \gamma^{i_k}(x_k - \varepsilon(x_k)), \quad i_1 + \cdots + i_k \geq i.
01190: \]
01191: L'anneau gradue associe a cette filtration sera note $\Gr^{\bullet}(X)$, et c'est lui
01192: qui remplacera dans le Seminaire l'anneau de Chow de $X$. On peut d'ailleurs
01193: montrer que dans le cas envisage ici l'homomorphisme canonique
01194: \begin{equation}
01195: \Gr^{\bullet}(X) \longrightarrow \Gr^*_{\mathrm{top}}(X)
01196: \end{equation}
01197: est un isomorphisme modulo torsion, ce qui justifie le point de vue suivant
01198: lequel cet anneau tensorise par $\mathbf{Q}$ peut jouer en tous points le meme role
01199: que l'anneau de Chow tensorise par $\mathbf{Q}$. D'autre part, on a fait exactement
01200: ce qu'il fallait pour que la formule (2.3) (ou on remplace $\Fil_{\mathrm{top}}$ par $\Fil_{\lambda}$)
01201: definisse une theorie des classes de Chern a valeurs dans $\Gr^{\bullet}(X)$, dont les
01202: classes de Chern dans $\Gr^*_{\mathrm{top}}(X) \otimes \mathbf{Q}$ se deduisent a l'aide de l'homomorphisme ci-dessus.
01203: On montre (Exp. V), grace au fait (Exp. VI) que $K(X)$ est un $\lambda$-anneau special
01204: au sens de [RRR] (i.e. un $\lambda$-anneau au sens de Exp. V) que cette notion de clas-
01205: se de Chern satisfait a toutes les proprietes formelles habituelles, passees
01206: en revue dans [RRR].
01207: 
01208: On notera un avantage tres appreciable de la definition de $\Gr^{\bullet}(X)$ sur
01209: celle de $A(X)$ et de $\Gr^*_{\mathrm{top}}(X)$, c'est qu'elle garde un sens pour un schema $X$ abso-
01210: lument quelconque, et plus generalement pour un topos localement annele
01211: quelconque $X$, en prenant alors pour $K(X)$ le groupe de classes construit avec
01212: les Modules localement libres sur $X$, qui est encore un $\lambda$-anneau augmente,
01213: donc muni d'une filtration correspondante, permettant de definir un gradue
01214: associe $\Gr^{\bullet}(X)$. Comme prix de cette generalite, il faut cependant remarquer
01215: que lorsqu'on n'est pas dispose a negliger les phenomenes de torsion i.e. a
01216: tensoriser par $\mathbf{Q}$, les proprietes de $\Gr^{\bullet}(X)$ sont moins satisfaisantes a divers
01217: ergards que celles de ses deux concurrents (Exp. XIV nos 4). C'est ainsi que,
01218: supposant a nouveau que $f: X \to Y$ est un morphisme propre de schemas noetheriens
01219: reguliers munis de Modules inversibles amples, il n'est pas vrai en general
01220: que l'homomorphisme $f_{*}: K(X) \to K(Y)$ applique $\Fil^i K(X)$ dans $\Fil^{i-d} K(Y)$, ou $d$
01221: est la dimension relative de $X$ sur $Y$, et definisse par suite par passage aux
01222: gradues associes un homomorphisme de degre $-d$ de $\Gr^{\bullet}(X)$ dans $\Gr^{\bullet}(Y)$. Cepen-
01223: dant, on peut prouver (Exp. VII) qu'il en est ainsi a condition de tensoriser
01224: par $\mathbf{Q}$, de sorte qu'on obtient un homomorphisme de degre $-d$
01225: \begin{equation}
01226: f_{*}: \Gr^{\bullet}(X) \otimes \mathbf{Q} \longrightarrow \Gr^{\bullet}(Y) \otimes \mathbf{Q},
01227: \end{equation}
01228: jouant le role de l'homomorphisme ``image directe de cycles'' $A(X) \to A(Y)$.
01229: 
01230: Ceci pose, nous voyons donc que nous sommes en possession de tous les
01231: elements de structure pour donner un sens a la formule (\ref{eq:RR2}) lorsqu'on ne
01232: postule pas l'existence d'un corps de base, $f: X \to Y$ etant morphisme propre
01233: de schemas noetheriens reguliers admettant des Modules inversibles amples.
01234: 
01235: \subsection*{3. Le cas des schemas generaux}
01236: 
01237: Montrons maintenant comment on peut encore donner un sens a la formule
01238: (\ref{eq:RR2}) lorsqu'on abandonne l'hypothese de regularite faite sur $X, Y$, en la
01239: remplacant par une hypothese de nature locale sur le morphisme $f$, que nous
01240: allons preciser.
01241: 
01242: \subsubsection*{3.1. Distinction entre $K_{\bullet}(X)$ et $K^{\bullet}(X)$}
01243: 
01244: Une premiere difficulte provient du fait que l'observation qui avait
01245: servi de point de depart a la theorie [RRR], savoir que le groupe $K(X)$ des
01246: classes de faisceaux sur $X$ est le meme, qu'on le definisse en termes de fais-
01247: ceaux coherents, ou en termes de faisceaux localement libres, est liee de
01248: facon essentielle a l'hypothese de regularite sur $X$. Dans le cas ou $X$ est un
01249: schema (disons noetherien, pour simplifier) general, il convient de distin-
01250: guer entre ces deux groupes, que nous noterons respectivement $K_{\bullet}(X)$ et $K^{\bullet}(X)$,
01251: le premier ayant un caractere covariant en $X$ pour les morphismes propres, le
01252: deuxieme un caractere contravariant en $X$ pour les morphismes quelconques. On
01253: notera que $K^{\bullet}(X)$ est muni de facon naturelle d'une structure d'anneau (et
01254: meme de $\lambda$-anneau augmente), alors que sur $K_{\bullet}(X)$ il n'y a plus en general une
01255: structure d'anneau, mais seulement une structure de module sur $K^{\bullet}(X)$. L'homo-
01256: morphisme $x \mapsto x \cdot \Cl(\mathcal{O}_X)$ (ou $\Cl_{\bullet}(F)$ designe la classe de $F$ dans $K_{\bullet}(X)$) de
01257: $K^{\bullet}(X)$ dans $K_{\bullet}(X)$ n'est plus en general un isomorphisme. Pour notre formula-
01258: tion du theoreme de Riemann-Roch generalise, nous travaillerons exclusivement
01259: avec $K^{\bullet}(X)$, $K^{\bullet}(Y)$, a l'exclusion de $K_{\bullet}(X)$, $K_{\bullet}(Y)$. Comme nous avons signale
01260: dans 2.3., ces $\lambda$-anneaux augmentes, munis de leur $\lambda$-filtration naturelle,
01261: donnent naissance a des anneaux gradues $\Gr^{\bullet}(X)$, $\Gr^{\bullet}(Y)$, et des homomorphismes
01262: ``classes de Chern''
01263: \[
01264: c_i: K^{\bullet}(X) \longrightarrow \Gr^{\bullet}(X),
01265: \]
01266: et de meme pour $Y$. La difficulte vient cependant dans la definition d'un
01267: homomorphisme image directe
01268: \begin{equation}
01269: f_{*}: K^{\bullet}(X) \longrightarrow K^{\bullet}(Y),
01270: \end{equation}
01271: qui ne peut plus etre defini par la formule naive
01272: \[
01273: f_{*}(\Cl(F)) = \sum (-1)^i \Cl(R^i f_{*}(F)),
01274: \]
01275: car meme si $F$ est localement libre, les $R^i f_{*}(F)$ ne sont plus en
01276: general localement libres, ni meme de tor-dimension finie, et la formule
01277: precedente perd son sens. Une solution de cette difficulte est fournie par
01278: le point de vue des categories derivees. En effet, alors que les $R^i f_{*}(F)$
01279: sont en general des ``mauvais'' faisceaux, de tor-dimension infinie disons,
01280: auxquels on ne saurait associer de classe dans $K^{\bullet}(Y)$, cependant le complexe
01281: $\mathbf{R}f_{*}(F)$ sur $Y$, image directe de $F$ au sens des categories derivees, dont les
01282: $R^i f_{*}(F)$ sont des faisceaux de cohomologie, a tendance a etre excellent. De
01283: facon precise, supposons que $F$ soit de tor-dimension finie sur $\mathcal{O}_Y$, ce qui
01284: sera le cas si $f$ est de tor-dimension finie sur $\mathcal{O}_Y$, et si de plus $F$ est
01285: localement libre. Alors on montre (Exp. III) que le complexe $\mathbf{R}f_{*}(F)$ est
01286: ``parfait'', par quoi on entend ici: a faisceaux de cohomologie coherents
01287: (ce qu'on savait deja de toutes facons) et de tor-dimension globale finie,
01288: ou encore: sur tout ouvert affine $U$ de $Y$, ce complexe est isomorphe, dans
01289: la categorie derivee $\mathbf{D}(U)$, a un complexe a degres bornes et localement libre
01290: de type fini en chaque degre. Du point de vue de l'Algebre Homologique, les
01291: complexes parfaits sont ``aussi bons'' que les Modules localement libres, dont
01292: ils constituent la generalisation naturelle. Ils forment d'ailleurs une
01293: ``sous-categorie triangulee'' de la categorie derivee $\mathbf{D}(Y)$, i.e. une sous-
01294: categorie stable par les cones (ou ``mapping cylinders'') de morphismes. D'au-
01295: tre part, a toute categorie triangulee $C$ on associe un groupe $K(C)$ de classes,
01296: variante naturelle de la construction du groupe des classes d'une categorie
01297: abelienne, en prenant le groupe libre engendre par les objets de $C$, divise
01298: par les relations
01299: \[
01300: \Cl(K) - \Cl(K') - \Cl(K'')
01301: \]
01302: pour tout triangle distingue
01303: \[
01304: K' \longrightarrow K \longrightarrow K'' \longrightarrow K'[1].
01305: \]
01306: Ce qui remplace alors l'observation fondamentale de [RRR] dans le cas actuel
01307: est le fait que, lorsque le schema $Y$ admet un Module inversible ample, tout
01308: complexe parfait $L_{\bullet}$ sur $Y$ est globalement (et non seulement localement)
01309: isomorphe dans $\mathbf{D}(Y)$ a un complexe $L'_{\bullet}$ a degres bornes et localement libre de
01310: type fini en chaque degre (Exp. II); d'ou il resulte facilement (Exp. IV) que
01311: l'homomorphisme canonique
01312: \[
01313: K^{\bullet}(X) \longrightarrow K(\mathrm{Parf}(Y)),
01314: \]
01315: ou $\mathrm{Parf}(Y)$ est la categorie triangulee des complexes parfaits sur $Y$, est un
01316: isomorphisme. Donc a tout complexe parfait $L_{\bullet}$ sur $Y$ on peut associer une
01317: classe
01318: \[
01319: \Cl(L_{\bullet}) \in K^{\bullet}(Y),
01320: \]
01321: qui n'est d'ailleurs autre que
01322: \[
01323: \Cl(L_{\bullet}) = \sum (-1)^i \Cl(L'_i),
01324: \]
01325: ou $L'_{\bullet}$ est comme ci-dessus, de sorte que les $L'_i$ sont bien des Modules loca-
01326: lement libres sur $Y$ dont on peut prendre les classes dans $K^{\bullet}(Y)$.
01327: 
01328: Ceci pose, on definit un homomorphisme (3.1) chaque fois que $f$
01329: est un morphisme propre de tor-dimension finie de schemas noetheriens tels
01330: que $Y$ admette un Module inversible ample, en posant simplement
01331: \begin{equation}
01332: f_{*}(\Cl(F)) = \Cl(\mathbf{R}f_{*}(F))
01333: \end{equation}
01334: pour tout Module localement libre $F$ sur $X$ (de sorte que $\mathbf{R}f_{*}(F)$ est bien
01335: un complexe parfait sur $Y$, et le deuxieme membre est defini). Lorsque $X$
01336: egalement admet un Module inversible ample, alors $K^{\bullet}(X)$ peut s'interpreter
01337: egalement comme $K(\mathrm{Parf}(X))$, et l'homomorphisme (3.1) n'est autre que celui
01338: qui est associe, par fonctorialite, au foncteur exact de categories trian-
01339: gulees induit par
01340: \[
01341: \mathbf{R}f_{*}: \mathrm{Parf}(X) \longrightarrow \mathrm{Parf}(Y).
01342: \]
01343: 
01344: \subsubsection*{4.1. Le groupe $K^{\bullet}(X)$ des complexes parfaits}
01345: 
01346: Notons d'abord que pour un schema $X$ (ou plus generalement, un topos
01347: localement annele) quelconque, il semble hors de doute que le ``bon'' invariant
01348: qui doit remplacer l'anneau $K(X)$ de [RRR] n'est pas l'anneau $K^{\bullet}(X)$ considere
01349: precedemment, defini en termes de Modules localement libres sur $X$, mais bien
01350: le groupe $K(\mathrm{Parf}(X))$, envisage dans 3.1., defini en termes de complexes
01351: parfaits sur $X$, groupe que nous noterons ici $K^{\bullet}(X)$, et que nous distingu-
01352: erons du groupe ``naif'' de 3.1, que nous notons $K^{\bullet}_{\text{naif}}(X)$. Ce point de vue est
01353: impose par la necessite de definir un homomorphisme (3.1) image directe pour
01354: un morphisme $f: X \to Y$ propre et de tor-dimension finie, de schemas noetheriens
01355: disons (le cas non noetherien exigeant des proprietes de finitude locales
01356: un peu plus fortes sur $f$, cf. Exp. III). Avec la definition de $K^{\bullet}(X)$ adoptee
01357: maintenant, cet homomorphisme $f_{*}$ se definit en effet de facon triviale par
01358: la formule (3.2), ou maintenant $F$ est un complexe parfait quelconque sur $X$,
01359: ce qui implique que $\mathbf{R}f_{*}(F)$ est egalement parfait.
01360: 
01361: \subsubsection*{4.2. La structure de $\lambda$-anneau sur $K^{\bullet}(X)$}
01362: 
01363: Cette nouvelle definition de $K^{\bullet}(X)$ souleve cependant un nouveau pro-
01364: bleme d'algebre homologique, qui est celui d'une definition d'une $\lambda$-struc-
01365: ture sur $K^{\bullet}(X)$, compatible avec sa structure d'anneau naturelle provenant du
01366: produit tensoriel dans la categorie derivee. (NB Si $L_{\bullet}$ et $L'_{\bullet}$ sont des com-
01367: plexes parfaits, il en est de meme de leur produit tensoriel total $L_{\bullet} \otimes^{\mathbf{L}} L'_{\bullet}$).
01368: Plus precisement, il y aurait lieu de definir des operations ``puissances
01369: exterieures'' dans la categorie derivee $\mathbf{D}(X)$, transformant complexes parfaits
01370: en complexes parfaits. Cette question n'a pas encore ete tiree au clair a
01371: l'heure actuelle, et comme il a ete dit dans l'Introduction, elle constitue
01372: la lacune la plus serieuse dans le Seminaire. Une fois cette question resolue,
01373: et procedant comme dans 2.3, on trouve une filtration naturelle sur $K^{\bullet}(X)$,
01374: permettant de former l'anneau gradue associe $\Gr^{\bullet}(X)$, qui jouera le role de
01375: l'anneau de Chow, et de definir une theorie des classes de Chern
01376: \[
01377: c_i: K^{\bullet}(X) \longrightarrow \Gr^{\bullet}(X).
01378: \]
01379: 
01380: \subsubsection*{4.3. Hypotheses supplementaires}
01381: 
01382: Pour donner un sens a la formule (\ref{eq:RR2}) pour un morphisme propre loca-
01383: lement d'intersection complete $f: X \to Y$ de schemas noetheriens, il faut encore
01384: inclure dans l'enonce de cette formule les inclusions non triviales (3.4)
01385: (ou $d$ est la dimension relative virtuelle) pour permettre de definir (3.3),
01386: et de definir enfin la classe tangente relative virtuelle
01387: \[
01388: T_f \in K^{\bullet}(X).
01389: \]
01390: Lorsque $f$ admet une factorisation (2.1) par une immersion dans un
01391: schema relatif lisse, la definition de $T_f$ n'offre pas de nouvelle difficul-
01392: te. On peut d'ailleurs la reformuler dans ce cas, d'une facon plus conforme
01393: a l'esprit ``categories derivees'', en considerant l'homomorphisme canonique
01394: bien connu
01395: \begin{equation}
01396: N_{X/X'} \longrightarrow \Omega^1_{X'/Y} \otimes_{\mathcal{O}_{X'}} \mathcal{O}_X
01397: \end{equation}
01398: induit par la differentielle exterieure
01399: \[
01400: d_{X'/Y}: \mathcal{O}_{X'} \longrightarrow \Omega^1_{X'/Y},
01401: \]
01402: et le complexe de chaines
01403: \begin{equation}
01404: L_{X/Y} \in \mathbf{D}(X)
01405: \end{equation}
01406: defini grace a la formule ci-dessus ``en prolongeant par des zeros'', complexe egal en degre
01407: 1 resp. 0 a la source resp. au but de l'homomorphisme ci-dessus. La formule (2.1)
01408: prend alors la forme plus sympathique
01409: \begin{equation}
01410: T_f = -\sum (-1)^i \Cl(L_i) = \Cl(L_{X/Y}^{\vee}),
01411: \end{equation}
01412: formule qui a un sens dans $K^{\bullet}(X)$ lorsque $f$ est un morphisme localement d'in-
01413: tersection complete, puisque dans ce cas le complexe $L_{X/Y}$ est manifestement
01414: parfait. L'assertion d'invariance de $T_f$ relativement a la factorisation choi-
01415: sie de $f$ peut alors se preciser facilement en une assertion d'invariance
01416: pour le complexe $L_{X/Y}$: modulo isomorphisme canonique dans la categorie deri-
01417: vee $\mathbf{D}(X)$, ce complexe est egalement independant de la factorisation choisie
01418: (Exp. VIII).
01419: 
01420: \subsubsection*{4.4. Construction du complexe cotangent relatif}
01421: 
01422: Lorsque $f: X \to Y$ est un morphisme quelconque de schemas, on peut enco-
01423: re construire un objet canonique $L_{X/Y}$ de la categorie derivee $\mathbf{D}(X)$, defini
01424: a isomorphisme unique pres, qui sur chaque ouvert affine de $X$ au dessus d'un
01425: ouvert affine de $Y$ soit isomorphe au complexe construit suivant le principe
01426: expose a l'alinea precedent (ou il convient cependant dans le cas general de
01427: remplacer ``lisse'' par ``formellement lisse'', de sorte que toute algebre $B$ sur
01428: un anneau $A$ apparait comme quotient d'une algebre formellement lisse, par
01429: exemple une algebre de polynomes). Pour une telle definition il n'est pas
01430: possible de proceder par simple ``recollement'' a partir des constructions
01431: affines, car les objets de la categorie derivee $\mathbf{D}(X)$ sont de nature essen-
01432: tiellement non recollables. Voici une construction generale de $L_{X/Y}$, vala-
01433: ble plus generalement pour tout morphisme $f: X \to Y$ de topos (commutative-
01434: ment) anneles: Soit $A = f^{*}(\mathcal{O}_Y)$, de sorte que $A$ est un Anneau sur l'es-
01435: pac $X$, et $B$ une $A$-Algebre. Soit $T \to B$ un homomorphisme de faisceaux d'en-
01436: sembles qui engendre $B$ comme $A$-Algebre, i.e. tel que l'homomorphisme de
01437: $A$-Algebres
01438: \[
01439: \phi: A[T] \longrightarrow B
01440: \]
01441: correspondant soit un epimorphisme de faisceaux d'ensembles, ou $A[T]$ designe
01442: la $A$-Algebre commutative libre engendree par $T$. Soit $J$ le noyau de $\phi$, et considerons comme dans 4.3 le complexe de chaines de longueur 1 defini par l'homomorphisme naturelle
01443: \[
01444: J/J^2 \longrightarrow \Omega^1_{A[T]/A} \otimes_{A[T]} B.
01445: \]
01446: Il n'est pas difficile de voir que, a isomorphisme canonique pres dans la
01447: categorie $\mathbf{D}(X)$, le complexe ainsi construit est independant du choix de
01448: $T$ et de $T \to B$, de sorte qu'il merite la designation de \textbf{complexe cotangent
01449: relatif} de $X$ sur $Y$. On verifie egalement que pour tout ouvert affine
01450: de $X$ au dessus d'un ouvert affine de $Y$, il est donne (a isomorphisme canoni-
01451: que pres) par le procede de 4.3.
01452: 
01453: Dans le cas ou $f$ est localement d'intersection complete, on voit de
01454: plus que l'on obtient meme un complexe parfait, ce qui permet alors de defi-
01455: nir encore un element $T_f$ dans $K^{\bullet}(X)$ par la formule (4.3). Signalons qu'ici
01456: encore, il n'aurait pas ete possible en general de definir $T_f$ comme un ele-
01457: ment de $K^{\bullet}_{\text{naif}}(X)$, et rien ne permet de supposer que $T_f$ soit globalement
01458: isomorphe dans $\mathbf{D}(X)$ a un complexe borne a composantes localement libres de
01459: type fini. Cela est donc une justification supplementaire du point de vue
01460: avance dans 4.1 en faveur du $K^{\bullet}(X)$ ``sophistique'', de preference a $K^{\bullet}_{\text{naif}}(X)$.
01461: 
01462: \subsubsection*{4.5. Conclusion}
01463: 
01464: Sous reserve de la question soulevee dans 4.2, nous avons donc enco-
01465: re reuni tous les elements de structure pour donner un sens a (\ref{eq:RR2}) pour
01466: un morphisme propre localement d'intersection complete de schemas quelcon-
01467: ques (l'hypothese noetherienne faite implicitement dans les sections qui
01468: precedent pouvant en fait s'eliminer sans difficulte essentielle). Cela ne
01469: signifie pas malheureusement que, meme sous cette reserve, la formule en
01470: question soit etablie, meme dans le cas ou $Y$ est le spectre d'un corps
01471: algebrique clos de caracteristique $p > 0$, et que $X$ est un schema algebrique
01472: propre et lisse de dimension 3 disons. Voir commentaires dans Exp. XIV no 2.
01473: 
01474: \subsection*{5. Variante analytique}
01475: 
01476: Sous la forme presentee dans le no 4, la formule de Riemann-Roch (\ref{eq:RR2})
01477: a egalement un sens pour un morphisme propre localement d'intersection com-
01478: plete d'espaces analytiques complexes (ou d'espaces rigide-analytiques au
01479: sens de Tate). Mais bien sur la demonstration donnee dans [RRR] ne s'appli-
01480: que a ce cas que lorsqu'on suppose de plus que le morphisme $f$ est projectif,
01481: et comme dans le cas algebrique, il semble qu'il faille une idee essentiel-
01482: lement nouvelle pour traiter le cas general. La formule que nous envisageons
01483: ici est une formule a valeurs dans un ``anneau de Chow analytique'' $\Gr^{\bullet}(X)$ ou
01484: plutot $\Gr^{\bullet}(X) \otimes \mathbf{Q}$, construit suivant le principe explique dans 4.2 (la diffi-
01485: culte qui y est rencontree s'evanouissant d'ailleurs du fait qu'on est en
01486: caracteristique nulle, comme on le signale dans Exp. XIV no 1).
01487: 
01488: Lorsque $X$ et $Y$ sont non singuliers, il est possible de donner une autre
01489: version de la formule de Riemann-Roch en procedant comme suit. Soit d'abord
01490: $X$ un espace analytique complexe quelconque, $X^{\text{top}}$ le meme espace muni du fais-
01491: ceau des fonctions complexes continues. On a donc un morphisme canonique
01492: d'espaces anneles
01493: \[
01494: \varepsilon: X^{\text{top}} \longrightarrow X,
01495: \]
01496: qui definit, grace au caractere contravariant evident du foncteur $K(X)$ en
01497: l'espace annele variable $X$, un homomorphisme d'anneaux
01498: \begin{equation}
01499: K^{\bullet}(X) \longrightarrow K^{\bullet}(X^{\text{top}}).
01500: \end{equation}
01501: Supposons pour simplifier $X^{\text{top}}$ compact; alors, comme dans le cas (deja si-
01502: gnale dans 3.1) d'un schema admettant un module inversible ample, on voit
01503: (Exp. II) que tout complexe parfait sur $X^{\text{top}}$ est globalement isomorphe a un
01504: complexe borne localement libre en chaque degre, i.e. a un complexe de fi-
01505: bres vectoriels complexes sur $X^{\text{top}}$, et que par suite $K^{\bullet}(X^{\text{top}})$ est isomorphe
01506: canoniquement au groupe $K(X^{\text{top}})$ bien connu des topologues, defini en
01507: termes de fibres vectoriels complexes sur l'espace compact $X^{\text{top}}$. On peut
01508: donc interpreter l'homomorphisme (5.1) comme etant un homomorphisme du $K^{\bullet}(X)$
01509: defini en termes de la structure analytique, dans l'anneau bien connu $K(X^{\text{top}})$
01510: des topologues. Si maintenant $f: X \to Y$ est un morphisme d'espaces analyti-
01511: ques non singuliers compacts, on sait lui associer grace a Atiyah-Hirzebruch
01512: 
01513: 
01514: un homomorphisme de groupes
01515: \begin{equation}
01516: f^{\text{top}}_{!}: K^{\bullet}(X^{\text{top}}) \longrightarrow K^{\bullet}(Y^{\text{top}}).
01517: \end{equation}
01518: Utilisant de meme l'homomorphisme $f_{!}$ de (3.1) defini par la formule (3.2)
01519: (qui utilise ici implicitement le theoreme de finitude de Grauert [4] pour
01520: un morphisme propre d'espaces analytiques) on trouve un carre d'homomor-
01521: phismes naturels, ne faisant intervenir que des groupes du type $K^{\bullet}$ (a l'ex-
01522: clusion d'anneaux du type anneaux de Chow ou de cohomologie):
01523: \begin{equation}
01524: \begin{tikzcd}
01525: K^{\bullet}(X) \arrow[r] \arrow[d, "f_{!}"'] & K^{\bullet}(X^{\text{top}}) \arrow[d, "f^{\text{top}}_{!}"] \\
01526: K^{\bullet}(Y) \arrow[r] & K^{\bullet}(Y^{\text{top}})
01527: \end{tikzcd}
01528: \end{equation}
01529: 
01530: La formule de type Riemann-Roch a laquelle nous faisions allusion con-
01531: siste a affirmer la commutativite de ce carre. On notera la difference de
01532: nature entre cette assertion (qui ne neglige pas les phenomenes de torsion,
01533: et se place dans un groupe $K^{\bullet}(Y^{\text{top}})$ de nature essentiellement discrete), et
01534: la formule de Riemann-Roch du type envisage au no 4, se placant dans un
01535: groupe du type ``Chow'' qui n'est plus de nature discrete, mais qu'on doit
01536: tensoriser en revanche par $\mathbf{Q}$. Notons que la formule (5.3) est demontree
01537: en tous cas lorsque $X$ et $Y$ sont des varietes projectives non singulieres
01538: [7], et comme la formule du type envisage au no 4, qui est demontree dans ce
01539: cas dans [RRR]. Il s'agit la d'une variante plausible du theoreme de l'index
01540: de Atiyah-Singer (et de ses generalisations dues a Shih [6] et Atiyah), qui
01541: appelle evidemment une generalisation commune. Une fois prouvee la commuta-
01542: tivite de (5.3), on en conclurait aussitot, a l'aide du ``theoreme de Riemann-Roch
01543: differentiable'' de Atiyah-Hirzebruch [1], une variante cohomologique de
01544: la formule de Riemann-Roch analytique complexe, savoir la commutativite du
01545: carre
01546: \begin{equation}
01547: \begin{tikzcd}
01548: K^{\bullet}(X) \arrow[r, "\ch"] \arrow[d, "f_{!}"'] & H^{2*}(X, \mathbf{Q}) \arrow[d, "f_{*}"'] \\
01549: K^{\bullet}(Y) \arrow[r, "\ch"'] & H^{2*}(Y, \mathbf{Q})
01550: \end{tikzcd}
01551: \end{equation}
01552: ou les fleches horizontales sont deduites des classes de Chern cohomologi-
01553: ques, deduites de l'homomorphisme (5.1) et de l'homomorphisme classe de
01554: Chern ordinaire
01555: \[
01556: K^{\bullet}(X^{\text{top}}) \longrightarrow H^{2*}(X, \mathbf{Z}),
01557: \]
01558: la premiere fleche verticale est la meme que dans (5.3), et la deuxieme est
01559: l'homomorphisme de Gysin en cohomologie rationnelle.
01560: 
01561: \subsection*{6. L'invariant covariant $K_{\bullet}(X)$}
01562: 
01563: Dans tout ceci, il n'a ete guere question que de l'invariant contrava-
01564: riant $K^{\bullet}(X)$, a l'exclusion de l'invariant covariant $K_{\bullet}(X)$ dont l'existence
01565: a ete signalee dans 3.1. Si on considere $K^{\bullet}(X)$ comme donnant un analogue de
01566: l'anneau de cohomologie entiere d'un espace topologique, on doit considerer
01567: $K_{\bullet}(X)$ comme donnant un analogue de l'homologie entiere, la structure de
01568: Module de $K_{\bullet}(X)$ sur $K^{\bullet}(X)$ correspondant alors a l'operation de cap-produit.
01569: Le fait que lorsque $X$ est regulier, l'homomorphisme naturel $K^{\bullet}(X) \to K_{\bullet}(X)$
01570: soit un isomorphisme, doit etre regarde comme l'analogue du theoreme de dua-
01571: lite de Poincare sur les varietes topologiques orientees, affirmant que le
01572: cap-produit avec la classe d'homologie fondamentale definit un isomorphisme
01573: de la cohomologie sur l'homologie. Ces analogies suggerent d'ailleurs l'op-
01574: portunite de definir egalement, pour un polyedre fini $X$ mettons, un inva-
01575: riant $K_{\bullet}(X)$ de nature covariante, dont les relations a l'homologie entiere
01576: (suite spectrale, homomorphisme de Chern...) seraient analogues a celles
01577: existant entre $K^{\bullet}(X)$ et la cohomologie entiere; et il ne semble pas non
01578: plus exclu qu'il puisse exister, en termes de ces groupes $K_{\bullet}(X)$, des varian-
01579: tes du theoreme de Riemann-Roch differentiable, valable pour des espaces
01580: plus generaux que des varietes differentiables compactes.
01581: 
01582: Revenant au cas ou $X$ est un schema noetherien quelconque, il convient
01583: de munir $K_{\bullet}(X)$ d'une filtration croissante, $\operatorname{Fil}_i K_{\bullet}(X)$ etant engendre par
01584: des classes $\operatorname{Cl}_{\bullet}(F)$, avec $F$ faisceau coherent sur $X$ tel que $\dim \operatorname{supp} F \leq i$.
01585: On prouvera (Exp. X) la compatibilite de cette filtration avec la structure
01586: de module et la filtration de $K^{\bullet}(X)$, i.e. la formule
01587: \begin{equation}
01588: \Fil^i K^{\bullet}(X) \cdot \Fil_j K_{\bullet}(X) \subset \Fil_{j-i} K_{\bullet}(X),
01589: \end{equation}
01590: ce qui permet de munir le gradue associe au module filtre $K_{\bullet}(X)$, note
01591: \[
01592: \Gr_{\bullet}(X),
01593: \]
01594: d'une structure de module sur $\Gr^{\bullet}(X)$. Ses elements peuvent d'ail-
01595: leurs s'interpreter comme des classes de cycles algebriques sur $X$, (pour une
01596: relation d'equivalence moins fine que l'equivalence rationnelle, lorsque $X$
01597: est de type fini sur un corps de base $k$).
01598: 
01599: Cette analogie n'est d'ailleurs pas purement verbale, et peut se preci-
01600: ser, apres passage aux gradues $\Gr^{\bullet}(X)$, $\Gr_{\bullet}(X)$ associes, par une relation de
01601: compatibilite explicite entre la structure de module sur $\Gr_{\bullet}(X)$ et l'opera-
01602: tion ``cap''.
01603: 
01604: Si $f: X \to Y$ est un morphisme propre de schemas noetheriens, alors
01605: \begin{equation}
01606: f_{*}: K_{\bullet}(X) \longrightarrow K_{\bullet}(Y)
01607: \end{equation}
01608: est compatible avec les filtrations, et definit un homomorphisme de groupes
01609: gradues
01610: \begin{equation}
01611: f_{*}: \Gr_{\bullet}(X) \longrightarrow \Gr_{\bullet}(Y),
01612: \end{equation}
01613: donnant lieu a une ``formule de projection'' i.e. qui est $\Gr^{\bullet}(Y)$-lineaire,
01614: formule qui se reduit par passage aux gradues de la formule de projection
01615: analogue pour l'homomorphisme (6.2) (qui est $K^{\bullet}(Y)$-lineaire).
01616: 
01617: Dans le point de vue propose dans le Seminaire, une theorie des in-
01618: tersections maniable, sur les schemas noetheriens $X$ pas necessairement re-
01619: guliers, doit consister en l'etude combinee des invariants contravariants
01620: $K^{\bullet}(X)$, $\Gr^{\bullet}(X)$ et des invariants covariants $K_{\bullet}(X)$, $\Gr_{\bullet}(X)$, qui jouent des
01621: roles de nature essentiellement differente, et dont les proprietes se com-
01622: pl\`etent mutuellement. C'est ainsi que les groupes covariants $K_{\bullet}(X)$, $\Gr_{\bullet}(X)$
01623: se calculent mieux pour certaines fibrations non propres, grace a la ``suite
01624: exacte d'homotopie'' qu'on peut etablir pour eux ([RRR] et Exp. IX); par
01625: exemple on prouve (Exp. IX) des isomorphismes canoniques tels que
01626: \[
01627: K_{\bullet}(X[t]) \simeq K_{\bullet}(X), \quad \Gr_{\bullet}(X[t]) \simeq \Gr_{\bullet}(X),
01628: \]
01629: qui tombent en defaut en general pour $K^{\bullet}$ lorsque $X$ est un schema non regu-
01630: ier. En revanche, $K^{\bullet}(X)$ et $\Gr^{\bullet}(X)$ se pretent bien aux calculs grace a leur
01631: structure d'anneaux.
01632: 
01633: Lorsqu'on considere des schemas $X$ propres sur un corps, la projection
01634: \[
01635: \pi_X: X \longrightarrow \text{(point)}
01636: \]
01637: definit un homomorphisme ``caracteristique d'Euler-Poincare''
01638: \begin{equation}
01639: \chi = \pi_{X*}: K_{\bullet}(X) \longrightarrow K_{\bullet}(\text{point}) = \mathbf{Z},
01640: \end{equation}
01641: induisant sur le sous-groupe $\Gr_0(X) = \Fil_0 K_{\bullet}(X)$ l'homomorphisme ``degre des
01642: 0-cycles''
01643: \begin{equation}
01644: \deg: \Gr_0(X) \longrightarrow \mathbf{Z}.
01645: \end{equation}
01646: L'homomorphisme de multiplication
01647: \[
01648: \Gr^i(X) \times \Gr_i(X) \longrightarrow \Gr_0(X),
01649: \]
01650: compose avec l'homomorphisme degre (6.5), fournit alors un accouplement
01651: \begin{equation}
01652: \Gr^i(X) \times \Gr_i(X) \longrightarrow \mathbf{Z},
01653: \end{equation}
01654: qui explique qu'a certains egards $\Gr^{\bullet}(X)$ et $\Gr_{\bullet}(X)$ jouent des roles duels,
01655: tout comme la cohomologie et l'homologie. Ainsi, si $f: X \to Y$ est un mor-
01656: phisme de schemas algebriques propres, les deux homomorphismes
01657: \[
01658: f^{*}: \Gr^{\bullet}(Y) \longrightarrow \Gr^{\bullet}(X) \quad \text{et} \quad f_{*}: \Gr_{\bullet}(X) \longrightarrow \Gr_{\bullet}(Y)
01659: \]
01660: formellement transposes l'un de l'autre, au sens des accouplements
01661: (6.6).
01662: 
01663: Les homomorphismes (6.4) et (6.5) et l'accouplement (6.6) qu'on en
01664: deduit, peuvent etre consideres comme la source des relations numeriques
01665: en theorie des intersections sur les schemas algebriques propres sur un
01666: corps donne $k$. Ce point de vue sera developpe avec quelques details dans
01667: Exp. X, independant pour l'essentiel des Exp. VI, VII, VIII (centres sur la
01668: demonstration du theoreme de Riemann-Roch) et de Exp. IX. On y donnera ega-
01669: lement quelques notions sur le comportement de $K^{\bullet}(X)$, $\Gr^{\bullet}(X)$ par specialisa-
01670: tion.
01671: 
01672: \subsection*{7. Le foncteur determinant}
01673: 
01674: Comme dans toute theorie d'intersection, il importe de preciser la
01675: place que prend dans la theorie le groupe de Picard (ou groupe des classes
01676: de diviseurs dans la terminologie classique). Moyennant une legere restric-
01677: tion, automatiquement satisfaite lorsque $X$ est par exemple quasi-compact,
01678: on etablit dans Exp. X un isomorphisme canonique
01679: \begin{equation}
01680: c_1: \Pic(X) \simeq \Gr^1(X).
01681: \end{equation}
01682: qui montre le role exact du groupe de Picard dans la theorie que nous pro-
01683: posons. Ce role justifie une etude plus detaillee de ce groupe de Picard
01684: dans les exposes XII, XIII, qui sont logiquement independants du texte
01685: du Seminaire. Une etude de l'equivalence numerique et de certains theoremes
01686: de finitude pour le groupe de Picard (annonces sans demonstration dans [6])
01687: est faite par Kleiman dans Exp. XIII, en utilisant des techniques de nature
01688: purement numerique (n'utilisant pas la theorie du foncteur de Picard et de
01689: sa representabilite), dues a lui-meme et a Mumford, nettement plus simples
01690: que les demonstrations initiales de Grothendieck. Le seul resultat refrac-
01691: taire aux methodes de Kleiman (et dont il est oblige de faire usage) est
01692: [EGA III, 7.7], dont la demonstration et certains raffinements (incluant
01693: certains theoremes de representabilite pour le foncteur de Picard) occupe
01694: l'expose XII. Ce dernier releve d'ailleurs d'un esprit et d'une technique
01695: entierement etrangers au reste du Seminaire, se rattachant a EGA $IV_3$; cet ex-
01696: pose est logiquement independant de tous les autres.
01697: 
01698: Enfin, dans Exp. XI, de nature nettement plus generale que XII et XIII
01699: et plus proche de l'esprit general du Seminaire, on etudie, sur un topos
01700: localement annele quelconque, un foncteur remarquable appele ``foncteur deter-
01701: minant''
01702: \begin{equation}
01703: \detn: \mathrm{Parf}(X) \longrightarrow \Inv(X)
01704: \end{equation}
01705: de la categorie des complexes parfaits sur $X$, les morphismes etant les seuls
01706: isomorphismes au sens de la categorie derivee $\mathbf{D}(X)$, dans la categorie des
01707: Modules inversibles sur $X$. Ce foncteur est decrit a peu de choses pres par
01708: l'exigence d'etre de ``nature locale'', et d'etre donne, pour un complexe $L_{\bullet}$
01709: bourne et a composantes localement libres, par la formule
01710: \begin{equation}
01711: \detn(L_{\bullet}) = \bigotimes_i \det(L_i)^{(-1)^i},
01712: \end{equation}
01713: ou $\det(L_i)$ designe la puissance exterieure de degre maximum du Module locale-
01714: ment libre $L_i$. Bien que cet expose aussi soit logiquement independant du res-
01715: te du Seminaire a l'exception de l'expose I, il constitue neanmoins un element
01716: important du ``yoga'' general propose dans le Seminaire, et developpe pour l'es-
01717: sentiel dans les exposes I a XI.
01718: 
01719: 
01720: \section*{Bibliographie}
01721: \addcontentsline{toc}{section}{Bibliographie}
01722: 
01723: \begin{enumerate}[label={[\arabic*]},leftmargin=*]
01724: \item M. Atiyah et F. Hirzebruch, Vector bundles and homogeneous spaces,
01725: \textit{Symposia in Pure Mathematics}, vol. 3, Differential Geometry,
01726: p. 7--38, 1961.
01727: 
01728: \item A. Borel et J.P. Serre, Le theoreme de Riemann-Roch, \textit{Bull. Soc. Math.}
01729: \textit{France}, t. 86, p. 97--136 (1958).
01730: 
01731: \item H. Grauert, Ein Theorem der analytischen Garbentheorie, \textit{Pub. Math.}
01732: I.H.E.S, no 5, p. 5--64 (1960).
01733: 
01734: \item A. Grothendieck, Classes de faisceaux et theoreme de Riemann-Roch,
01735: notes mimeographiees, Princeton 1957, (cite [RRR] et reproduit
01736: en Appendice a l'Exp. 0 du Seminaire).
01737: 
01738: \item A. Grothendieck, Technique de descente et theoremes d'existence en
01739: Geometrie Algebrique, dans \textit{Fondements de la Geometrie Algebrique},
01740: Extraits du Seminaire Bourbaki 1957--1962, Secretariat Mathema-
01741: tique, Paris.
01742: 
01743: \item W. Shih, Seminaire de Topologie a l'IHES, 1966/67.
01744: 
01745: \item M. Atiyah et F. Hirzebruch, The Riemann-Roch theorem for analytic
01746: embeddings, \textit{Topology}, vol. 1, p. 151--166 (1962).
01747: \end{enumerate}
01748: 
01749: 
01750: \section*{Appendice}
01751: \addcontentsline{toc}{section}{Appendice -- Classes de Faisceaux et Theoreme de Riemann-Roch}
01752: 
01753: \begin{center}
01754: \Large\textbf{CLASSES DE FAISCEAUX ET THEOREME DE RIEMANN-ROCH}\\[0.5cm]
01755: \textit{par A. Grothendieck}\footnote{Ceci est la reproduction textuelle du rapport cite dans l'Introduction
01756: du present Seminaire. Les notes de bas de page indiquees par des signes
01757: (*), (**) ont ete rajoutees en Octobre 1967.}\\[0.3cm]
01758: \end{center}
01759: 
01760: \begin{center}
01761: \textbf{CHAPITRE I}
01762: 
01763: \medskip
01764: \textbf{$\lambda$-Anneaux (Preliminaires formels)}
01765: 
01766: \medskip
01767: \begin{tabular}{lp{8cm}r}
01768: $\S$ 1. & Definitions & \dotfill 1 \\
01769: $\S$ 2. & Exemples & \dotfill 4 \\
01770: $\S$ 3. & Le $\lambda$-anneau defini par un anneau gradue & \dotfill 8 \\
01771: $\S$ 4. & Les operations $\rho^p(N, x)$ & \dotfill 13 \\
01772: \end{tabular}
01773: \end{center}
01774: 
01775: 
01776: \begin{center}
01777: \textbf{CHAPITRE II}
01778: 
01779: \medskip
01780: \textbf{Classes de faisceaux algebriques coherents et classes de Chern}
01781: 
01782: \medskip
01783: \begin{tabular}{lp{8cm}r}
01784: $\S$ 1. & La theorie de Chow & \dotfill 18 \\
01785: $\S$ 2. & Definition des classes de Chern des faisceaux algebriques coherents & \dotfill 22 \\
01786: $\S$ 3. & Generalites fonctorielles sur $K(X)$ & \dotfill 27 \\
01787: $\S$ 4. & Quelques resultats techniques & \dotfill 34 \\
01788: $\S$ 5. & Definition faisceautique des classes de Chern. Application a l'etude des morphismes d'injection & \dotfill 43 \\
01789: $\S$ 6. & Le theoreme de Riemann-Roch & \dotfill 49 \\
01790: \end{tabular}
01791: \end{center}
01792: 
01793: \begin{center}
01794: \textit{Demonstration du theoreme de Riemann-Roch (1er Novembre 1957).}
01795: \end{center}
01796: 
01797: 
01798: \section*{Chapitre I}
01799: \addcontentsline{toc}{subsection}{Chapitre I -- $\lambda$-Anneaux}
01800: 
01801: \begin{center}
01802: \Large\textbf{$\lambda$-ANNEAUX (PRELIMINAIRES FORMELS)}
01803: \end{center}
01804: 
01805: \subsection*{\S 1. Definitions}
01806: 
01807: On appelle \textbf{$\lambda$-anneau}\footnote{Dans ce seminaire, nous disons plutot pre-$\lambda$-anneau (V 2.1), reser-
01808: vant le nom de $\lambda$-anneau a ceux satisfaisant la condition supplemen-
01809: taire de la page 3 (cf. V 2.4).} un anneau commutatif $K$ avec unite, muni d'une
01810: famille d'applications
01811: \[
01812: \lambda^i: K \longrightarrow K \qquad (i \text{ entier } \geq 0)
01813: \]
01814: satisfaisant aux conditions suivantes:
01815: \begin{equation}
01816: \lambda^0(x) = 1, \quad \lambda^1(x) = x, \quad \lambda^n(x+y) = \sum_{i=0}^{n} \lambda^i(x) \lambda^{n-i}(y).
01817: \end{equation}
01818: 
01819: Posons, pour tout $x \in K$:
01820: \begin{equation}
01821: \lambda_t(x) = \sum_{n \geq 0} \lambda^n(x) t^n \in K[[t]];
01822: \end{equation}
01823: les relations precedentes expriment alors que $x \mapsto \lambda_t(x)$ est un homomor-
01824: phisme du groupe additif $K$ dans le groupe multiplicatif $1 + K[[t]]^{+}$ des
01825: series formelles dans $K$, d'augmentation 1, relevant l'homomorphisme cano-
01826: nique $x \mapsto 1 + x \cdot t$.
01827: 
01828: Ainsi, sur l'anneau $\mathbf{Z}$ existe une unique $\lambda$-structure compatible avec
01829: sa structure d'anneau et telle que
01830: \[
01831: \lambda_t(1) = 1 + t.
01832: \]
01833: On a alors
01834: \begin{equation}
01835: \lambda_t(n) = (1+t)^n, \quad \text{d'ou } \lambda^i(n) = \binom{n}{i}.
01836: \end{equation}
01837: La notion de $\lambda$-homomorphisme de $\lambda$-anneaux est claire; un $\lambda$-anneau
01838: augmente est un $\lambda$-anneau muni d'un $\lambda$-homomorphisme dans $\mathbf{Z}$.
01839: 
01840: Pour definir la notion de $\lambda$-anneau special, les considerations sui-
01841: vantes sont utiles. Soit $k$ un anneau commutatif avec unite, soit
01842: $\widetilde{K} = 1 + k[[t]]^{+}$ le groupe des series formelles a coefficients dans $k$,
01843: d'augmentation 1; on va y introduire une loi de $\lambda$-anneau dont la structure
01844: additive soit la multiplication des series formelles. La multiplication
01845: de cet anneau sera notee $f \circ g$. Elle est entierement determinee par la
01846: condition que les coefficients $c_n(f \circ g)$ de $f \circ g$ soient donnes
01847: par des polynomes universels a coefficients entiers en les coefficients
01848: $a_i = c_i(f)$ et $b_j = c_j(g)$ de $f$ et $g$ respectivement:
01849: \begin{equation}
01850: c_n = P_n(a_1, \ldots, a_n; b_1, \ldots, b_n),
01851: \end{equation}
01852: qu'elle soit distributive par rapport a l'``addition'' $f \cdot g$, et soit donnee
01853: pour des facteurs lineaires par:
01854: \begin{equation}
01855: (1 + at) \circ (1 + bt) = 1 + abt.
01856: \end{equation}
01857: 
01858: Les $P_n$ se calculent par des calculs de fonctions symetriques; $P_n$ est
01859: isobare de poids $n$ en les $a_i$ et isobare de poids $n$ en les $b_j$ ($a_i$ et $b_j$
01860: etant de poids $i$) et d'ailleurs symetrique en $a = (a_i)$ et $b = (b_j)$.
01861: 
01862: De meme, les $\lambda^i(f)$ sont determines sans ambiguite par la condition que
01863: les coefficients $c_i(\lambda^n(f)) = c_i^n$ de $\lambda^n(f)$ se calculent au moyen de poly-
01864: nomes universels a coefficients entiers en les coefficients $c_i(f) = a_i$
01865: de $f$:
01866: \begin{equation}
01867: c_i^n = P_{i,n}(a_1, a_2, \ldots, a_{in}),
01868: \end{equation}
01869: que les relations (1.2) soient verifiees, et que, pour $i > 1$, on ait:
01870: \begin{equation}
01871: \lambda^i(1 + at) = 1.
01872: \end{equation}
01873: (1 est l'element nul de $1 + k[[t]]^{+}$). Les $P_{i,n}$ se calculent
01874: encore par des calculs de fonctions symetriques, et l'on constate que
01875: $P_{i,n}$ est un polynome isobare de poids $in$ en les $a_i$. Muni des ope-
01876: rations precedentes, $\widetilde{K}$ devient un $\lambda$-anneau, son element nul est 1, son element
01877: unite est $1+t$. Notons aussi la relation:
01878: \begin{equation}
01879: (1 + at) \circ f = f(at).
01880: \end{equation}
01881: 
01882: Soit maintenant $K$ un $\lambda$-anneau quelconque. On dit que $K$ est \textbf{special}
01883: si l'homomorphisme additif de $K$ dans $1 + K[[t]]^{+}$ est un $\lambda$-homomorphisme.
01884: Cela signifie donc qu'on a les formules:
01885: \begin{equation}
01886: \lambda_t(1) = 1 + t, \quad \lambda_t(x \circ y) = \lambda_t(x) \circ \lambda_t(y), \quad \lambda_t(\lambda^n(x)) = \lambda^n(\lambda_t(x)),
01887: \end{equation}
01888: ou encore, explicitement:
01889: \begin{equation}
01890: \begin{aligned}
01891: \lambda^i(1) &= 0 \text{ si } i \geq 1 \\
01892: \lambda^n(x \cdot y) &= P_n(\lambda^1 x, \ldots, \lambda^n x; \lambda^1 y, \ldots, \lambda^n y) \\
01893: \lambda^m(\lambda^n(x)) &= P_{m,n}(\lambda^1 x, \ldots, \lambda^{mn} x)
01894: \end{aligned}
01895: \end{equation}
01896: ou les $P_n$ et les $P_{m,n}$ sont les polynomes universels envisages plus haut.
01897: 
01898: Il en resulte en particulier les formules (1.5). (On peut montrer que
01899: le $\lambda$-anneau $1 + k[[t]]^{+}$ construit avec un anneau commutatif $k$ est spe-
01900: cial; mais nous n'aurons pas besoin de ce fait).
01901: 
01902: \subsection*{\S 2. Exemples}
01903: 
01904: \textbf{1)} Soient $G$ un groupe et $k$ un anneau commutatif; on considere le
01905: groupe $K(G)$, quotient du groupe libre engendre par les classes de repre-
01906: sentations de $G$ dans des $k$-modules de type fini par le sous-groupe engen-
01907: dre par les elements de la forme $(\rho \oplus \rho') - (\rho) - (\rho')$ (on note $\oplus$ la
01908: somme directe). Le produit tensoriel et les puissances exterieures defi-
01909: nissent dans $K(G)$ une structure de $\lambda$-anneau. On peut aussi passer au
01910: quotient par les elements de la forme $(\rho) - (\rho') - (\rho'')$ ou $\rho$ est une
01911: presentation extension de $\rho'$ par $\rho''$, triviale comme extension de $k$-modules.
01912: On obtient le groupe $K_{\bullet}(G)$ des classes reduites de representations de $G$,
01913: et l'on constate que la loi de $\lambda$-anneau passe au quotient.
01914: 
01915: Lorsque $k$ est un corps algebriquement clos de caracteristique 0\footnote{Il est inutile que $k$ soit algebriquement clos, cf. (VI 3.3) pour un
01916: resultat encore plus general (englobant tous ceux qui seront utilises
01917: dans le present rapport).},
01918: on peut montrer que $K(G)$, et a fortiori $K_{\bullet}(G)$, est un $\lambda$-anneau special;
01919: il n'en est plus de meme si la caracteristique est $\neq 0$, mais cependant
01920: si $k$ est algebriquement clos\footnotemark[1], $K_{\bullet}(G)$ est encore un $\lambda$-anneau special.
01921: \addtocounter{footnote}{1}
01922: \footnotetext{Condition inutile, cf. note de bas de page precedente.}
01923: 
01924: Pour le prouver, on est ramene au cas ou $G$ est un produit $\mathrm{GL}(m,k) \times \mathrm{GL}(n,k)$
01925: et ou l'on ne considere que les representations rationnelles de $G$;
01926: notre assertion resulte de la determination explicite de $K_{\bullet}(G)$ qui
01927: sera faite plus bas (exemple 3); enfin, lorsque $k$ est de caracteristique
01928: 0, la complete reductibilite des representations rationnelles de $G$ montre
01929: que $K(G) = K_{\bullet}(G)$.
01930: 
01931: \medskip
01932: \textbf{Variantes} de l'exemple precedent: $G$ est un groupe algebrique defini
01933: sur $k$ et on ne considere que les representations rationnelles de $G$ de-
01934: finies sur $k$; on peut considerer le cas d'un groupe topologique, ou
01935: de Lie, et des representations continues, etc\ldots\footnote{On peut construire d'autres exemples de $\lambda$-anneaux, par exemple l'en-
01936: semble des classes d'equivalence de formes quadratiques, l'ensemble
01937: des classes de representations d'une algebre de Lie, etc\ldots\ Nous pro-
01938: mettons que nous n'aurons pas besoin de tous ces exemples pour traiter
01939: le theoreme de Riemann-Roch !}
01940: 
01941: \medskip
01942: \textbf{2)} Soit $X$ une variete algebrique, irreductible ou non, definie sur
01943: le corps $k$. On pourra considerer le groupe quotient du groupe libre en-
01944: gendre par les classes de fibres vectoriels sur $X$ (definis sur $k$), par
01945: le sous-groupe engendre par les elements de la forme $(E \oplus E') - (E) - (E')$,
01946: sur lequel les operations de $\lambda$-anneau seront definies par le produit
01947: tensoriel et les puissances exterieures. Si $X$ est complete, c'est un
01948: groupe libre engendre par les fibres vectoriels indecomposables sur $X$
01949: (grace a Krull-Remak-Schmidt); la connaissance de cet anneau et de sa
01950: base equivaut a la classification complete des fibres vectoriels sur
01951: $X$ et la connaissance des operations tensorielles elementaires sur ces
01952: fibres. Si $k$ est algebriquement clos de caracteristique 0, toutes les
01953: autres operations tensorielles sur des classes de fibres sont alors con-
01954: nues, a l'aide de ce qui va suivre dans 3). On peut faire un passage au
01955: quotient plus draconien, par le groupe engendre par les elements
01956: $(E) - (E') - (E'')$ ou $E$ est un fibre vectoriel extension de $E'$ par $E''$.
01957: On trouve alors un $\lambda$-anneau, note $K(\mathcal{E}(X))$, ou $\mathcal{E}(X)$ designe la categorie
01958: additive des fibres vectoriels sur $X$, definis sur $k$. En utilisant les
01959: resultats de l'exemple 1), on montre que si $k$ est algebriquement clos\footnotemark[2], le $\lambda$-anneau $K(\mathcal{E}(X))$ est special, et qu'il en est de meme du
01960: premier $\lambda$-anneau defini ici, si la caracteristique de $k$ est 0, mais non
01961: dans le cas general.
01962: \addtocounter{footnote}{1}
01963: \footnotetext{On peut construire d'autres exemples de $\lambda$-anneaux, par exemple l'en-
01964: semble des classes d'equivalence de formes quadratiques, l'ensemble
01965: des classes de representations d'une algebre de Lie, etc\ldots}
01966: 
01967: \medskip
01968: \textbf{3) Determination de $K(G)$ et de $K_{\bullet}(G)$ dans certains cas importants:}
01969: On suppose le corps $k$ algebriquement clos, le groupe $G$ algebrique affine
01970: connexe et ``globalement reductif'' (le radical de $G$ est isomorphe a un
01971: tore) et l'on ne considere que les representations lineaires rationnelles.
01972: En caracteristique 0, la complete reductibilite des representations line-
01973: aires rationnelles de $G$ montre que $K(G) = K_{\bullet}(G)$ est le groupe libre
01974: engendre par les classes des representations rationnelles irreductibles.
01975: 
01976: Or, si $T$ est un tore maximal de $G$, on a $K(T) = K_{\bullet}(T) = \mathbf{Z}[\hat{T}]$, algebre
01977: du groupe $\hat{T}$ des caracteres rationnels de $T$; on a $\hat{T} \simeq \mathbf{Z}^r$ si $r$ est le
01978: rang de $G$ (ceci est independant de la caracteristique). Le groupe de
01979: Weyl $W$ opere dans $\hat{T}$, donc dans $K(T)$, d'une maniere qui ne depend que des
01980: operations de $W$ dans le groupe libre $\hat{T}$. On a un homomorphisme evident
01981: de $\lambda$-anneaux, deduit de l'injection de $T$ dans $G$:
01982: \begin{equation}
01983: K_{\bullet}(G) \longrightarrow K_{\bullet}(T) = K(T)^W.
01984: \end{equation}
01985: ou $K(T)^W$ denote l'ensemble des points fixes de $W$ dans $K(T)$. La theorie
01986: des representations lineaires de $G$\footnote{Cf. Seminaire Chevalley 1956/58 Groupes de Lie algebriques (Ecole
01987: Normale Superieure).} demontre alors le theoreme suivant:
01988: 
01989: \begin{theorem}
01990: L'homomorphisme (1.13) est un isomorphisme.
01991: \end{theorem}
01992: 
01993: On voit donc que $K_{\bullet}(G)$, une fois connu le ``type'' de $G$ (caracterise
01994: par $W$ operant sur un groupe $\mathbf{Z}^r$), ne depend pas de la caracteristique,
01995: et que c'est un $\lambda$-anneau special. On a une base canonique de $K_{\bullet}(G)$
01996: correspondant aux orbites de $W$ dans $\hat{T}$; une autre base s'obtient ainsi:
01997: soient $(\rho_i)_{1 \leq i \leq r}$ les classes de representation de $G$ correspondant aux
01998: representations fondamentales de $G/\mathrm{rad}\, G$, et $(\sigma_j)_{1 \leq j \leq s}$ une base du
01999: groupe des homomorphismes rationnels de $G$ dans $k^{*}$; on a donc
02000: \[
02001: r + s = r' \quad \text{(rang de } G\text{)}.
02002: \]
02003: 
02004: \begin{corollary}
02005: L'anneau $K_{\bullet}(G)$ s'identifie au sous-anneau
02006: \[
02007: \mathbf{Z}[\rho_1, \ldots, \rho_r; \sigma_1, \sigma_1^{-1}, \ldots, \sigma_s, \sigma_s^{-1}]
02008: \]
02009: du corps des fonctions rationnelles en les $\rho_i$ et les $\sigma_j$ a coefficients dans $\mathbf{Q}$.
02010: \end{corollary}
02011: 
02012: En particulier, $K_{\bullet}(G)$ n'a pas de diviseurs de 0; si $G = \prod \mathrm{GL}(n_i, k)$,
02013: l'anneau $K_{\bullet}(G)$ s'explicite ainsi: soit $\rho_{ij}$ la representation de $G$ deduite
02014: de la representation identique de $\mathrm{GL}(n_i, k)$; si l'on considere alors
02015: le corps des fonctions rationnelles sur $\mathbf{Q}$ en des indeterminees $(e_{ij})$
02016: ($1 \leq j \leq n_i$), alors $K_{\bullet}(G)$ est le sous-anneau engendre par ces inde-
02017: terminees et les inverses des $e_{ij}$. Cela justifie notre assertion
02018: que, du moins en caracteristique 0, les operations $\lambda^i$ permettent de
02019: reconstituer toutes les operations tensorielles sur les fibres vectoriels;
02020: ceci est encore vrai sur un corps algebriquement clos de caracteristique
02021: quelconque, pourvu qu'on prenne des classes reduites.
02022: 
02023: \medskip
02024: \textbf{Remarques:} \textbf{a)} Dans l'exemple 1), l'anneau $K(G)$ est augmente, l'aug-
02025: mentation etant definie par le degre des representations, et par passage
02026: au quotient, cette augmentation definit une augmentation sur $K_{\bullet}(G)$. De
02027: meme, les $\lambda$-anneaux envisages dans l'exemple 2) sont augmentes lorsque
02028: $X$ est $k$-irreductible; dans $K(\mathcal{E}(X))$, l'augmentation est definie par le
02029: rang d'un fibre vectoriel, c'est-a-dire la dimension de ses fibres.
02030: 
02031: \textbf{b)} Lorsque le corps $k$ n'est pas algebriquement clos, il est plausible
02032: que $K(G)$ est un $\lambda$-anneau special en caracteristique 0, et de meme pour
02033: $K_{\bullet}(G)$ en toutes caracteristiques.
02034: 
02035: \textbf{c)} En caracteristique non nulle, le theoreme 1.1 m'a ete gracieusement
02036: fourni par Chevalley.
02037: 
02038: 
02039: \subsection*{\S 3. Le $\lambda$-anneau defini par un anneau gradue}
02040: 
02041: Soit $A$ un anneau gradue commutatif avec unite; on supposera pour
02042: simplifier que les degres de $A$ sont positifs et que $A^0 = \mathbf{Z}$. On designe
02043: par $A^{+}$ le produit des $A^i$ ($i \geq 1$) (c'est un anneau commutatif avec unite), par
02044: $A_{+}$ le noyau de l'augmentation evidente $A^{+} \to A^0 = \mathbf{Z}$. Alors $1 + A_{+}$ est
02045: un sous-groupe du groupe multiplicatif des elements inversibles de $A$.
02046: 
02047: Considerons le groupe abelien
02048: \begin{equation}
02049: \mathbf{A} = \mathbf{Z} \times (1 + A_{+})
02050: \end{equation}
02051: dont les elements seront notes $[n, x]$ avec $n \in \mathbf{Z}$ ($x \in 1 + A_{+}$).
02052: Nous ecrirons ce groupe additivement, ainsi:
02053: \begin{equation}
02054: [n, x] + [n', x'] = [n + n', x \cdot x'];
02055: \end{equation}
02056: l'element nul de ce groupe est $[0, 1]$. Nous allons y introduire une struc-
02057: ture de $\lambda$-anneau augmente, compatible avec sa structure additive, l'aug-
02058: mentation etant $\varepsilon: [n, x] \mapsto n$. Le produit $[m, x] \cdot [n, y]$ est comple-
02059: tement caracterise par le fait qu'il s'ecrive $[mn, x^n y^m \cdot (x * y)]$, ou $x * y$
02060: est un element de $1 + A_{+}$ dont les composantes s'expriment comme polynomes
02061: universels a coefficients entiers en les $x^i$ et les $y^j$:
02062: \begin{equation}
02063: (x * y)^i = Q_i(x^1, \ldots, x^i; y^1, \ldots, y^i) \quad (i \geq 1)
02064: \end{equation}
02065: ou $Q_i$ est isobare de poids $i$. De plus, $x * y$ doit etre distri-
02066: butif par rapport a l'``addition'' $x \cdot y$ et pour des facteurs ``lineaires'' doit
02067: se reduire a ce qui suit:
02068: \begin{equation}
02069: [1, 1+x] \cdot [1, 1+y] = [1, 1+x+y].
02070: \end{equation}
02071: 
02072: Les polynomes $Q_i$ s'evaluent par des calculs de fonctions symetriques
02073: elementaires (il s'agit du calcul des classes de Chern d'un produit ten-
02074: soriel de fibres vectoriels a l'aide des classes de Chern $x^i$ et $y^j$ des
02075: facteurs et de leurs dimensions respectives $m$ et $n$). De (1.17) on deduit:
02076: \[
02077: Q_i(x^1, \ldots, x^i; y^1, \ldots, y^i) = x^i + y^i + \text{(termes utilisant les } x^p, y^q \text{ avec } p, q < i\text{)},
02078: \]
02079: et un produit $x * y$ quelconque se calcule terme a terme en decomposant
02080: ``formellement'' $x$ en un produit $\prod_{i=1}^{m} (1+a_i)$ et $y$ en un produit $\prod_{j=1}^{n} (1+b_j)$,
02081: ou les $a_i$ et les $b_j$ sont de degre 1.
02082: 
02083: De cette maniere, $\mathbf{A}$ devient un anneau commutatif augmente, avec
02084: une unite $[1, 1]$, et on peut l'interpreter comme l'anneau obtenu par
02085: adjonction d'une unite a l'anneau $1 + A_{+}$. De plus, sur $\mathbf{A}$, on definit
02086: une filtration en posant $\mathbf{A}^0 = \mathbf{A}$ et en notant $\mathbf{A}^n$ l'ensemble des $[0, x]$
02087: avec $x^i = 0$ pour $1 \leq i < n$; cette filtration est compatible avec le produit,
02088: car on a
02089: \begin{equation}
02090: [0, x] \in \mathbf{A}^m, \; [0, y] \in \mathbf{A}^n \implies [0, x] \cdot [0, y] - [0, 1] \in \mathbf{A}^{m+n},
02091: \end{equation}
02092: les termes non critiques etant de degre $> m+n$. Cette formule montre
02093: que l'anneau gradue $G(\mathbf{A})$ associe a $\mathbf{A}$ s'identifie a $A$ au point de vue
02094: additif, mais non au point de vue multiplicatif; on peut dire seulement
02095: que l'application $\eta: A \to G(\mathbf{A})$ definie par
02096: \begin{equation}
02097: \eta(a) = \mathrm{cl}([1, 1+a]) \in G^i(\mathbf{A}) \quad \text{pour } a \in A^i
02098: \end{equation}
02099: est un homomorphisme d'anneaux.
02100: 
02101: On introduit dans $\mathbf{A}$ des operations $\lambda^i$ verifiant les axiomes des
02102: $\lambda$-anneaux (\S 1, formules (1.2)), bien determinees par la condition que
02103: les composantes de $\lambda^i[n, x]$ se deduisent des composantes $x^j$ de $x$ par des
02104: polynomes universels a coefficients entiers
02105: \begin{equation}
02106: \lambda^i[n, x]^j = Q_{i,j,n}(x^1, x^2, \ldots, x^{ij})
02107: \end{equation}
02108: ($Q_{i,j,n}$ est isobare de poids $j$), et que l'on ait
02109: \begin{equation}
02110: \lambda^i[1, 1+x] = [0, 1] \text{ pour } i > 1.
02111: \end{equation}
02112: Les polynomes $Q_{i,j,n}$ s'evaluent par un calcul de fonctions symetriques ele-
02113: mentaires, qui est le calcul des classes de Chern des puissances exte-
02114: rieures d'un fibre vectoriel, a l'aide des classes de Chern et du rang
02115: de ce dernier. On constate alors que $\mathbf{A}$ est un $\lambda$-anneau special augmente
02116: et que les $\mathbf{A}^i$ sont stables par les operations $\lambda^j$.
02117: 
02118: Le $\lambda$-anneau augmente $\mathbf{A}$ peut etre considere comme un foncteur en $A$;
02119: par suite l'automorphisme principal $x = \sum_{i \geq 0} x^i \mapsto \bar{x} = \sum_{i \geq 0} (-1)^i x^i$
02120: de $A$ definit un automorphisme de $\mathbf{A}$, note par le meme symbole $\nabla$.
02121: 
02122: \begin{proposition}
02123: Pour $x \in \mathbf{A}$ de la forme $[n, 1 + x^1 + \cdots + x^n]$ (on a donc
02124: $x^i = 0$ pour $i > n$), on a $\lambda^i(x) = 0$ pour $i > n$, on a $\lambda^n(x) = [1, 1+x^n]$
02125: et enfin
02126: \begin{equation}
02127: \lambda_t(x) = \sum_{i=0}^{n} [1, 1+x^i] \cdot t^i \cdot \lambda_{-t}(\bar{x}^{n-i}).
02128: \end{equation}
02129: \end{proposition}
02130: 
02131: De cette proposition et du fait que les $\mathbf{A}^i$ sont stables par les
02132: operations $\lambda^j$, on deduit que la formule (1.22) est valable chaque fois
02133: que $\varepsilon(x) = n$, que l'on ait ou non $x^i = 0$ pour $i > n$.
02134: 
02135: \bigskip
02136: 
02137: La formule (1.22) admet les generalisations suivantes: soit $K$ un
02138: anneau quelconque et par ailleurs arbitraire; pour $x \in K$, on pose
02139: \begin{equation}
02140: \gamma^n(x) = (-1)^n \sum_{i=0}^{n} (-1)^i \binom{n-1+i}{n-i} \lambda^i(x) = \lambda^n(x+n-1),
02141: \end{equation}
02142: (l'identite des deux derniers membres resulte d'une recurrence immediate).
02143: 
02144: On a en particulier $\gamma^0(x) = 1$ et $\gamma^1(x) = x$; on pose ensuite:
02145: \begin{equation}
02146: \gamma_t(x) = \sum_{n \geq 0} \gamma^n(x) t^n \in K[[t]];
02147: \end{equation}
02148: d'ou l'on deduit la formule:
02149: \begin{equation}
02150: \gamma_t(x) = \lambda_{t/(1-t)}(x),
02151: \end{equation}
02152: ce qui equivaut a
02153: \begin{equation}
02154: \lambda_s(x) = \gamma_{s/(1+s)}(x).
02155: \end{equation}
02156: Ceci prouve que les $\lambda^i$ s'expriment a l'aide des $\gamma^j$ et que
02157: $\gamma_t(x+y) = \gamma_t(x) \cdot \gamma_t(y)$ et par suite que les applications $\gamma^i$ defini-
02158: ssent une nouvelle structure de $\lambda$-anneau sur $K$. La formule (1.22) montre
02159: alors que pour $K = \mathbf{A}$ et tout $x \in \mathbf{Z}$:
02160: \begin{equation}
02161: \gamma^n([0, 1+x] - \varepsilon([0, 1+x])) = [0, 1 + (n-1)! \, x^n + \ldots],
02162: \end{equation}
02163: et par suite, la composante $x^n$ s'obtient, au facteur $(n-1)!$ pres, en reduisant modulo $\mathbf{A}^{n+1}$ l'element $\gamma^n(x - \varepsilon(x))$ de $\mathbf{A}^n$.
02164: 
02165: \bigskip
02166: 
02167: \textbf{L'homomorphisme de Chern:}
02168: 
02169: Pour tout anneau gradue $A$ verifiant les conditions precedentes, on
02170: va definir un homomorphisme d'anneaux augmentes:
02171: \begin{equation}
02172: \ch: \mathbf{A} \longrightarrow A \otimes \mathbf{Q}.
02173: \end{equation}
02174: Cet homomorphisme est defini sans ambiguite par la condition d'etre un
02175: homomorphisme additif, de transformer 1 en 1, d'etre tel que les com-
02176: posantes de $\ch(x)$ en degre $> 0$ soient donnees par des polynomes univer-
02177: sels isobares a coefficients rationnels en les composantes de $x$, et que
02178: \begin{equation}
02179: \ch([1, 1+x^1]) = \exp(x^1 t) = \sum_{n=0}^{\infty} \frac{(x^1)^n}{n!} t^n.
02180: \end{equation}
02181: On obtient alors aussitot l'expression explicite de $\ch$: soit $\eta$ l'endomor-
02182: phisme additif de $A \otimes \mathbf{Q}$ defini pour $x$ de degre $i$ par
02183: \begin{equation}
02184: \eta(x^i) = (-1)^{i-1} (i-1)! \, x^i;
02185: \end{equation}
02186: on a alors
02187: \begin{equation}
02188: \ch([n, 1 + \sum_{i=1}^{N} x^i]) = n + \eta^{-1}\bigl(\log(1 + \sum_{i=1}^{N} \eta(x^i))\bigr).
02189: \end{equation}
02190: Comme on peut ``inverser'' cette formule, du moins si $A$ n'a qu'un nombre
02191: fini de composantes homogenes, on voit que dans ce cas, on a un isomor-
02192: phisme $\ch$ de l'anneau $\mathbf{A} \otimes \mathbf{Q}$ sur l'anneau $A \otimes \mathbf{Q}$ (ce qui elucide comple-
02193: tement la structure de $\mathbf{A} \otimes \mathbf{Q}$).
02194: 
02195: Rappelons que d'apres Hirzebruch [1], une serie formelle $f \in \mathbf{Q}[[t]]$ a
02196: coefficients dans $\mathbf{Q}$ definit par des formules polynomiales universelles un homomorphisme addi-
02197: tif $f_{*}: \mathbf{A} \otimes \mathbf{Q} \to \mathbf{A} \otimes \mathbf{Q}$; lorsque $f(t) = t/(1 - \exp(-t))$,
02198: on pose simplement $f_{*} = \Todd$ (l'initiale de Todd). On etend la defi-
02199: nition de $\Todd$ a $x \in K$. Ceci dit, on a le resultat suivant:
02200: 
02201: \begin{proposition}
02202: Soit $N \in \mathbf{A}$ de la forme $[q, 1+N^1+\cdots+N^q]$ (ona
02203: $N^i = 0$ pour $i > q$). Si l'on pose:
02204: \[
02205: \gamma_t(N-q) = \sum_{i=0}^{\infty} \gamma^i(N-q) t^i, \quad \lambda_t(N) = \sum_{i=0}^{q} \lambda^i(N) t^i,
02206: \]
02207: on aura:
02208: \begin{equation}
02209: \ch(\lambda_{-t}(N)) = (-1)^q \, \Todd(N) \cdot \ch(\gamma_t(N-q)).
02210: \end{equation}
02211: \end{proposition}
02212: 
02213: Ce qui equivaut a:
02214: \begin{equation}
02215: \ch(\lambda_{-1}(N)) = (-1)^q \, \Todd(N).
02216: \end{equation}
02217: 
02218: 
02219: \subsection*{\S 4. Les operations $\rho^p(N, x)$}
02220: 
02221: Soit $\mathbf{A}$ l'anneau des series formelles, a coefficients dans $\mathbf{Z}$, en
02222: deux suites infinies d'indeterminees, notees $r^i N$ et $\lambda^j x$ ($i, j \geq 1$). On pose
02223: $\lambda^0 N = N$ et $\lambda^0 x = x$, et l'on introduit sur $\mathbf{A}$ l'unique structure de $\lambda$-anneau
02224: special pour lequel on ait $\lambda^i(N) = r^i N$ et $\lambda^j(x) = \lambda^j x$ et pour lequel les
02225: operations $\lambda^i$ soient ``continues'' (si l'on se bornait aux polynomes en des
02226: indeterminees $r^i N$ et $\lambda^j x$, on aurait le ``$\lambda$-anneau libre engendre par $N$ et $x$'';
02227: notre anneau $\mathbf{A}$ en est un complete). On peut alors former
02228: \begin{equation}
02229: \lambda_{-t}(N) = \sum_{i=0}^{\infty} (-1)^i \lambda^i(N) t^i,
02230: \end{equation}
02231: qui est une serie formelle de terme constant 1, donc inversible; on
02232: peut donc definir sans ambiguite un element $\rho^p(N, x)$ de $\mathbf{A}$ en posant
02233: \begin{equation}
02234: \rho^p(N, x) \cdot \lambda_{-t}(N) = \lambda_t(\rho^p(N, x) \cdot \lambda_{-1}(N)).
02235: \end{equation}
02236: 
02237: \begin{theorem}
02238: Soit $\mathfrak{I}$ l'ideal ferme de $\mathbf{A}$ engendre par les $r^i N$ avec $i > q$;
02239: identifions $\mathbf{A}/\mathfrak{I}$ a l'anneau des series formelles en les $r^j x$ et les
02240: $r^k N$ ($1 \leq j \leq p$, $1 \leq k \leq q$). Alors $\rho^p(N, x)$ reduit modulo $\mathfrak{I}$ est un polynome, isobare de
02241: poids $p$ par rapport aux variables $r^j x$.
02242: \end{theorem}
02243: 
02244: Soient $N$ et $F$ deux espaces vectoriels de dimensions finies $q, q'$ sur
02245: un corps $k$ algebriquement clos de caracteristique 0; soit $n_0$ le noyau
02246: de l'application $(a_1, \ldots, a_q) \mapsto a_1 + \cdots + a_q$ de $N^q$ dans $N$. On
02247: fait operer le groupe symetrique $S_n$ sur $n_0^p$ donc sur $\Lambda^p n_0$ donc sur
02248: $\Lambda^p N \otimes \Lambda^p F$. Considerons cet espace
02249: vectoriel comme un module gradue sur $S_n \times \mathrm{GL}(N) \times \mathrm{GL}(F)$ et prenons
02250: sa ``partie alternee'' qui est un module gradue sur $\mathrm{GL}(N) \times \mathrm{GL}(F) = G$.
02251: On a alors
02252: \begin{equation}
02253: \rho^p(N, F) = \sum_i (-1)^i \Cl(\Lambda^i n_0^p \otimes \Lambda^p F).
02254: \end{equation}
02255: N.B. Si $M$ est un complexe de $G$-modules, on note $\sum (-1)^i \Cl(M_i)$ la somme al-
02256: ternee dans $K(G)$ des classes des composantes de $M$; c'est la caracteris-
02257: tique d'Euler-Poincare de $M$; dans la formule (1.33), le premier membre
02258: $\rho^p(N, F)$ designe l'element de l'anneau $K(G)$ obtenu en substituant aux va-
02259: riables $r^i N$ et $\lambda^j F$ les classes dans $K(G)$ des $G$-modules $\Lambda^i(N)$ et $\Lambda^j(F)$
02260: respectivement.
02261: 
02262: \textit{Demonstration:} En developpant le deuxieme membre de (1.32) par la
02263: deuxieme formule (1.11), on voit que la premiere assertion est prouvee
02264: si on prouve l'assertion analogue obtenue en faisant $x = 1$, i.e. si on
02265: prouve que dans le $\lambda$-anneau $\mathbf{Z}[r^1 N, \ldots, r^q N]$ ou $r^i N = 0$ pour $i > q$, l'ele-
02266: ment $\rho^p(N, 1)$ est divisible par $\lambda_{-1}(N)$. A priori, le quotient
02267: $\rho^p(N, 1) = \rho^p(N, 1) / \lambda_{-1}(N)$ est une serie formelle contenue dans le corps
02268: des quotients $K$ de l'anneau des series formelles en les $r^i N$ pour
02269: $1 \leq i \leq q$. D'autre part, designons aussi par $N$ un vectoriel de dimension
02270: $q$ sur un corps algebriquement clos $k$ de caracteristique 0, et consi-
02271: derons l'element $\rho^p(N) = \sum (-1)^i \Cl(\Lambda^i n_0^p) \in K(G)$, ou maintenant
02272: l'on a $G = \mathrm{GL}(N)$. On sait (\S 2, theoreme 1.1) que $K(G)$ s'identifie
02273: au sous-anneau de $K$ engendre par les $r^i N$ pour $1 \leq i \leq q$ et par l'inverse
02274: de $r^q N$. D'autre part, on verifie facilement dans cet anneau que l'on
02275: a $\rho^p(N) \cdot \lambda_{-1}(N) = \rho^p(N \cdot \lambda_{-1}(N))$. Comme on est dans un grand corps $K$ et que
02276: $r^q N \neq 0$, on en conclut $\rho^p(N) = \rho^p(N, 1)$; les deux membres sont dans
02277: l'intersection de $K(G)$ et de l'anneau des series formelles en les $r^i N$,
02278: donc sont des polynomes en les $r^i N$.
02279: 
02280: Pour prouver (1.33) dans le cas general, on verifie que le deuxieme
02281: membre satisfait dans $K(G)$ a la formule analogue a (1.32), c'est-a-dire
02282: que son produit par $\lambda_{-1}(N)$ est $\lambda_t(\rho^p(N \cdot \lambda_{-1}(N)))$. De meme, le produit de
02283: $\rho^p(N, F) \in K(G)$ par $\lambda_{-1}(N)$ est $\lambda_t(\rho^p(F \cdot \lambda_{-1}(N)))$, comme il resulte du fait
02284: que $K(G)$ est special. Comme $K(G)$ est integre (cf. \S 2) et que $\lambda_{-1}(N) \neq 0$
02285: (meme reference), on a bien demontre la formule (1.33).
02286: 
02287: \hfill \textit{c.q.f.d.}
02288: 
02289: \medskip
02290: \textbf{Remarques:} \textbf{a)} Lorsque la dimension $q'$ de $F$ est $\geq p$, alors (1.33) cara-
02291: c\'erise le polynome $\rho^p(N, x)$ en les $r^i N$ ($1 \leq i \leq q$) et les $\lambda^j x$ ($1 \leq j \leq p$).
02292: 
02293: \textbf{b)} Il est plausible que (1.33) est encore verifiee si le corps $k$
02294: n'est pas algebriquement clos\footnote{Cela se prouve en effet par la meme methode.}. Il faudrait resoudre cette question
02295: si l'on voulait prouver le theoreme de Riemann-Roch sur un corps de carac-
02296: teristique 0 non algebriquement clos par la methode qu'on va exposer. Le
02297: passage d'un $G$-module a sa partie alternee n'a d'ailleurs de sens raison-
02298: nable que si le corps de base est de caracteristique 0.
02299: 
02300: Soit maintenant $K$ un $\lambda$-anneau special quelconque et soit $N$ un ele-
02301: ment de $K$ tel que $r^i N = 0$ pour $i$ assez grand. On peut alors, pour tout
02302: $x \in K$ et tout entier $p \geq 1$, considerer les elements $\rho^p(N, x)$, qui s'ex-
02303: priment par des polynomes universels a coefficients entiers en les
02304: $r^i N$ et les $\lambda^j x$. On a $\rho^1(N, x) = x$ et l'on a la formule (1.32) dans $K$. Pour
02305: exprimer les proprietes des operations $x \mapsto \rho^p(N, x)$, il est commode
02306: d'introduire un nouveau $\lambda$-anneau note $K_N$. En tant qu'anneau commutatif,
02307: $K_N$ s'obtient par adjonction d'une unite a l'anneau dont le groupe addi-
02308: tif sous-jacent est $K$ et dont la multiplication est donnee par
02309: \begin{equation}
02310: x \cdot_N y = xy \cdot \rho^1(N, x) \cdot \rho^1(N, y).
02311: \end{equation}
02312: L'application $\lambda_t$ de $K_N$ dans $K_N[[t]]$ est definie par:
02313: \begin{equation}
02314: \lambda_t(x) = \sum_{p=0}^{\infty} \rho^p(N, x) t^p,
02315: \end{equation}
02316: ou l'on a pose $\rho^0(N, x) = [1, 0]$ (element unite de $K$). Je dis que,
02317: muni de ces operations, $K_N$ est un $\lambda$-anneau special. Il s'agit donc de
02318: verifier les formules $\lambda_t(X+Y) = \lambda_t(X) \lambda_t(Y)$, $\lambda_t(X \cdot_N Y) = \lambda_t(X) \circ \lambda_t(Y)$, et
02319: $\lambda_t(\lambda^n(X)) = \lambda^n(\lambda_t(X))$ pour $X, Y \in K_N$. On peut se borner a verifier ces
02320: formules lorsque $X$ et $Y$ sont dans $K$, et l'on constate qu'il suffit de
02321: le verifier lorsque $K$ est le $\lambda$-anneau special libre engendre par $X, Y$
02322: et $N$ soumis aux relations $r^i N = 0$ pour $i > q$. Mais l'homomorphisme de
02323: groupes
02324: \begin{equation}
02325: K_N \longrightarrow \mathbf{A}/\mathfrak{I},
02326: \end{equation}
02327: qui transforme unite en unite et sur $K$ se reduit a l'homomorphisme $x \mapsto xu$,
02328: est un homomorphisme compatible avec les structures d'anneaux et les
02329: applications $\lambda_t$, comme on le verifie aussitot. Dans le cas actuel,
02330: sa restriction a $K$ est injective; comme $K$ est un $\lambda$-anneau special, il
02331: en resulte facilement que les relations ci-dessus sont bien verifiees
02332: pour $X, Y \in K$, et ceci acheve de demontrer notre assertion.
02333: 
02334: D'apres la formule (1.23), il est naturel, pour $x \in K$ et $n \geq 1$, de
02335: poser
02336: \begin{equation}
02337: \gamma^n(N, x) = \sum_{i=0}^{n} \binom{n}{i} \rho^i(N, x),
02338: \end{equation}
02339: en d'autres termes, on pose $\gamma^n(N, x) = \rho^n(N, x_1)$ ou $x_1$ designe $x$ consi-
02340: dere comme element du $\lambda$-anneau $K_N$.
02341: 
02342: \begin{proposition}
02343: Soient $A$ un anneau gradue commutatif avec unite, $N$
02344: un element de $\mathbf{A}$ de la forme $N = [q, 1 + N^1 + \cdots + N^q]$ et $x$ un element
02345: quelconque de $\mathbf{A}$: $x = [n, 1 + \sum_{i \geq 1} x^i]$. Alors, pour tout entier $j \geq 0$,
02346: on a $\gamma^j(N, x) \in \mathbf{A}^j$, et la composante de degre $j$ de $\gamma^j(N, x)$ est
02347: de la forme
02348: \begin{equation}
02349: (\gamma^j(N, x))^j = C_{q,j}(N^1, \ldots, N^q; x^1, x^2, \ldots, x^j) \cdot x^j,
02350: \end{equation}
02351: ou $C_{q,j}$ est un polynome universel a coefficients entiers caracterise par
02352: la condition
02353: \begin{equation}
02354: C_{q,j}(N^1, \ldots, N^q; t, t^2, \ldots, t^j) = \binom{n+q+j-1}{j} t^j.
02355: \end{equation}
02356: \end{proposition}
02357: 
02358: \textit{Demonstration.} On peut supposer que $\mathbf{A}$ est l'anneau des polynomes a coef-
02359: ficients entiers en des indeterminees $x^j$ et $n^i$ (avec $1 \leq j \leq q$). Comme
02360: on a evidemment, combinant (1.32) et (1.37)
02361: \begin{equation}
02362: \gamma^p(N, x) \cdot \lambda_{-1}(N) = \gamma^p(N \cdot \lambda_{-1}(N)),
02363: \end{equation}
02364: et que, pour $n = q + j$, le deuxieme membre est de filtration $\geq q + j$
02365: (utiliser la formule (1.22 bis) en tenant compte de ce que $x \cdot \lambda_{-1}(N)$ est fini).
02366: 
02367: 
02368: \section*{Chapitre II}
02369: \addcontentsline{toc}{subsection}{Chapitre II -- Classes de faisceaux algebriques coherents et classes de Chern}
02370: 
02371: \begin{center}
02372: \Large\textbf{CLASSES DE FAISCEAUX ALGEBRIQUES COHERENTS ET CLASSES DE CHERN}
02373: \end{center}
02374: 
02375: \subsection*{\S 1. La theorie de Chow}
02376: 
02377: On supposera algebriquement clos, pour simplifier. Une grande partie de ce
02378: qui va suivre, et peut-etre tout, est cependant valable sans cette res-
02379: triction. Les espaces algebriques et les applications regulieres consideres
02380: seront definis sur $k$.
02381: 
02382: Un espace algebrique est dit quasi-projectif s'il est isomorphe a
02383: une partie localement fermee d'un espace projectif.
02384: 
02385: Soit $X$ une variete non singuliere quasi-projective connexe de dimen-
02386: sion $n$. On designe par $A(X)$ et l'on appelle \textbf{anneau de Chow} de $X$ l'anneau
02387: gradue des classes de cycles sur $X$ a equivalence rationnelle pres. $A^i(X)$
02388: etant forme des classes de cycles de dimension $n - i$. Le fait que $A(X)$
02389: soit bien un anneau gradue resulte du theoreme suivant de Chow:
02390: 
02391: \begin{theorem}[Chow]
02392: Tout cycle sur $X$ est rationnellement equivalent a un cycle
02393: non singulier. Etant donnes des elements de $A^i(X)$ et $A^j(X)$, on peut
02394: trouver des representants non singuliers de ces elements, soient
02395: $Y^i, Z^j$, qui se coupent partout transversalement.
02396: \end{theorem}
02397: 
02398: (Un cycle est dit non singulier si ses composantes sont des sous-varietes
02399: non singulieres. On dit que deux cycles se coupent partout transversalement
02400: si toute composante de l'un coupe partout transversalement toute compo-
02401: sante de l'autre, i.e. si les deux varietes secantes sont non singulieres
02402: aux points d'intersections et si leurs varietes tangentes y sont en
02403: position generale\ldots)
02404: 
02405: Soit $f$ un morphisme (i.e. une application reguliere) de $X$ dans $Y$
02406: ($X$ et $Y$ etant quasi-projectives et non singulieres), alors $f$ definit
02407: (en utilisant le theoreme precedent) un homomorphisme d'anneaux gradues:
02408: \begin{equation}
02409: f^{*}: A(Y) \longrightarrow A(X)
02410: \end{equation}
02411: (image reciproque de classes de cycles). Ainsi $A(X)$ devient un foncteur
02412: contravariant en $X$.
02413: 
02414: Soit $f$ un morphisme propre de $X$ dans $Y$ (i.e. dans la terminologie
02415: de Weil, le graphe de $f$ est complet au-dessus de $Y$, cf. [3]). Alors $f$
02416: definit aussi un homomorphisme additif
02417: \begin{equation}
02418: f_{*}: A(X) \longrightarrow A(Y)
02419: \end{equation}
02420: conservant la dimension des cycles, donc augmentant le degre de $m - n$
02421: (si $\dim X = n$, $\dim Y = m$). Ainsi, relativement aux morphismes propres, et faisant abstraction des
02422: structures multiplicatives et des graduations, $A(X)$ est un foncteur cova-
02423: riant en $X$. Notons la formule classique
02424: \begin{equation}
02425: f_{*}(x \cdot f^{*}(y)) = f_{*}(x) \cdot y.
02426: \end{equation}
02427: 
02428: Nous utiliserons aussi la theorie de Chow des classes de Chern. Elle
02429: consiste a attacher, a tout fibre vectoriel $E$ sur $X$ (toujours supposee
02430: quasi-projective non singuliere) des classes de Chern $c^i(E) \in A^i(X)$,
02431: ($i \geq 1$), de facon a satisfaire aux conditions suivantes; on pose $c^0(E) = 1$
02432: et
02433: \begin{equation}
02434: c_t(E) = \sum_{i=0}^{\infty} c^i(E) t^i \in A(X)[[t]].
02435: \end{equation}
02436: (voir chapitre I, \S 3 pour la definition de $A(X)$). Les conditions a veri-
02437: fier par une theorie des classes de Chern s'expriment alors ainsi:
02438: 
02439: \begin{equation}
02440: c(E) = c(E') \cdot c(E'')
02441: \end{equation}
02442: si le fibre vectoriel $E$ est extension de $E'$ par $E''$ (si l'on pose
02443: $c(E) = \sum c^i(E)$, cette condition s'exprime $c(E) = c(E') \cdot c(E'')$). La for-
02444: mule (2.5) permet d'etendre $c$ en un homomorphisme additif de $K(\mathcal{E}(X))$
02445: (cf. Chap. I, \S 2, exemple 2) dans $A(X)$:
02446: \begin{equation}
02447: c: K(\mathcal{E}(X)) \longrightarrow A(X),
02448: \end{equation}
02449: et l'on veut que $c$ soit un $\lambda$-homomorphisme de $\lambda$-anneaux; cette derniere
02450: condition s'ecrit aussi
02451: \begin{equation}
02452: c(E \otimes F) = c(E) \circ c(F), \quad c(\Lambda^i E) = \lambda^i(c(E)),
02453: \end{equation}
02454: pour deux fibres vectoriels $E$ et $F$ (bien entendu, $c$ est automatiquement
02455: compatible avec les augmentations). Soient $D$ un diviseur sur $X$ et $L(D)$
02456: le fibre vectoriel qu'il definit; on veut
02457: \begin{equation}
02458: c^1(L(D)) = \Cl(D) \in A^1(X)
02459: \end{equation}
02460: (ou pour un cycle quelconque $Z$, on note $\Cl(Z)$ sa classe dans $A(X)$). Enfin,
02461: on impose que les $c^i$ soient fonctoriels, i.e. que si $f$ est un morphisme
02462: de $X$ dans $Y$ et $E$ un fibre vectoriel sur $Y$, on ait
02463: \begin{equation}
02464: c^i(f^{*}(E)) = f^{*}(c^i(E)),
02465: \end{equation}
02466: ou $f^{*}(E)$ est l'image inverse de $E$ sur $X$. On peut aussi ecrire (2.9) sous
02467: la forme suivante: la notion d'image inverse de fibre definit un $\lambda$-homo-
02468: morphisme de $\lambda$-anneaux augmentes
02469: \begin{equation}
02470: f^{\circ}: K(Y) \longrightarrow K(X)
02471: \end{equation}
02472: (on ecrit ici $K(X)$ au lieu de $K(\mathcal{E}(X))$), et (2.9) signifie qu'on a commutativite
02473: dans le diagramme:
02474: \[
02475: \begin{tikzcd}
02476: K(Y) \arrow[r, "f^{\circ}"] \arrow[d, "c_Y"'] & K(X) \arrow[d, "c_X"] \\
02477: A(Y) \arrow[r, "f^{*}"'] & A(X)
02478: \end{tikzcd}
02479: \]
02480: (on a mis en indice au symbole $c$ l'espace sur lequel on le considere.)
02481: 
02482: \begin{theorem}
02483: Il existe une theorie des classes de Chern satisfaisant aux
02484: conditions (2.5), (2.7), (2.8) et (2.9).
02485: \end{theorem}
02486: 
02487: Nous admettrons aussi la proposition suivante:
02488: 
02489: \begin{proposition}
02490: Soit $Y$ une partie fermee de $X$ et soit $U$ son complemen-
02491: taire; si $x$ est un element de $A(X)$ qui induit $0$ sur $U$, il existe un
02492: representant de $x$ qui soit un cycle porte par $Y$.
02493: \end{proposition}
02494: 
02495: \begin{corollary}
02496: Soit $p = \dim X - \dim Y$. Alors, si $i < p$, l'homomorphisme
02497: $A^i(X) \to A^i(U)$ est bijectif. Tout element du noyau de $A^p(X) \to A^p(U)$
02498: est combinaisons lineaire de composantes de dimension $n - p$ de $Y$ (donc
02499: est proportionnel a $\Cl(Y)$ si $Y$ est irreductible).
02500: \end{corollary}
02501: 
02502: \subsection*{\S 2. Definition des classes de Chern des faisceaux algebriques coherents}
02503: 
02504: Soit $\mathcal{C}$ une categorie abelienne, et soit $\mathcal{C}_1$ une sous-classe de $\mathcal{C}$.
02505: On designe par $K(\mathcal{C}_1)$ le groupe quotient du groupe libre engendre par les
02506: classes (a isomorphisme pres) d'elements de $\mathcal{C}_1$ par le sous-groupe engen-
02507: dre par les elements de la forme $(E) - (E') - (E'')$, ou $E$ est une extension
02508: de $E'$ par $E''$ ($E, E', E'' \in \mathcal{C}_1$). On doit supposer que les classes d'elements de
02509: $\mathcal{C}_1$ pour la relation d'isomorphisme forment un ensemble. On designe par
02510: $\gamma_{\mathcal{C}_1}(E)$ la classe d'un $E \in \mathcal{C}_1$ dans $K(\mathcal{C}_1)$ et l'on omet l'indice $\mathcal{C}_1$ quand
02511: il n'y a pas de risque de confusion. Le groupe $K(\mathcal{C}_1)$ est solution d'un
02512: probleme universel relatif aux fonctions $E \mapsto f(E)$ definies sur $\mathcal{C}_1$ a
02513: valeurs dans un groupe abelien arbitraire $A$, et qui sont additives en ce
02514: sens que $f(E) = f(E') + f(E'')$ chaque fois que $E$ est extension de $E'$ par $E''$.
02515: 
02516: Soient $F$ un foncteur additif de $\mathcal{C}$ dans $\mathcal{C}'$ et $\mathcal{C}'_1$ une sous-classe de
02517: $\mathcal{C}'$ satisfaisant a la meme condition que $\mathcal{C}_1$ tels que pour tout $A \in \mathcal{C}_1$,
02518: les $R^qF(A)$ soient dans $\mathcal{C}'_1$ et nuls pour $q$ assez grand (on suppose que dans $\mathcal{C}$
02519: existent des resolutions injectives, de sorte que les foncteurs derives
02520: $R^qF$ de $F$ existent); alors $F$ definit un homomorphisme de $K(\mathcal{C}_1)$ dans $K(\mathcal{C}'_1)$,
02521: note encore $F$, par la formule
02522: \begin{equation}
02523: F(\gamma_{\mathcal{C}_1}(E)) = \sum_{i=0}^{\infty} (-1)^i \gamma_{\mathcal{C}'_1}(R^i F(E)).
02524: \end{equation}
02525: 
02526: Plus generalement, si on a un foncteur cohomologique $(F^i)$ de $\mathcal{C}$ dans $\mathcal{C}'$
02527: tel que pour tout $E \in \mathcal{C}$ les $F^i(E)$ soient dans $\mathcal{C}'_1$ et nuls sauf un nombre
02528: fini, on en deduit un homomorphisme de $K(\mathcal{C}_1)$ dans $K(\mathcal{C}'_1)$. Ces remarques
02529: s'etendent aux multifoncteurs; de plus, les suites spectrales ``classiques''
02530: donnent des proprietes de transitivite que nous n'expliciterons pas.
02531: En particulier, si $F$ est un foncteur exact, la formule (2.11) se
02532: simplifie
02533: \begin{equation}
02534: F(\gamma_{\mathcal{C}_1}(E)) = \gamma_{\mathcal{C}'_1}(F(E)).
02535: \end{equation}
02536: 
02537: Ceci s'applique en particulier si $\mathcal{C} = \mathcal{C}'$, si $F$ est le foncteur identique
02538: et $\mathcal{C}'_1 \subset \mathcal{C}_1$; on a alors un homomorphisme canonique de $K(\mathcal{C}'_1)$ dans $K(\mathcal{C}_1)$.
02539: 
02540: \begin{theorem}
02541: Soient $\mathcal{C}$ une categorie abelienne et $\mathcal{C}_1 \subset \mathcal{C}$, deux sous-classes
02542: telles que les classes (a isomorphisme pres) d'objets de $\mathcal{C}_1$ forment un
02543: ensemble. On suppose les conditions suivantes satisfaites:
02544: 
02545: \textit{(i)} Pour toute suite exacte $0 \to A' \to A \to A'' \to 0$ dans $\mathcal{C}$ telle
02546: que $A$ et $A''$ soient dans $\mathcal{C}_1$ (resp. $\mathcal{C}'_1$), on a $A' \in \mathcal{C}_1$ (resp. $A' \in \mathcal{C}'_1$).
02547: 
02548: \textit{(ii)} Soient $u: A \to B$ et $v: C \to B$ avec $A, B, C \in \mathcal{C}_1$ et $u$ surjectif;
02549: alors il existe $L \in \mathcal{C}_1$ et des homomorphismes $u': L \to C$ et $v': L \to A$
02550: tels que $u'$ soit surjectif et $vu' = uv'$.
02551: 
02552: \textit{(iii)} Pour tout $A \in \mathcal{C}_1$, il existe un entier $d = d(A)$ tel que pour
02553: toute resolution gauche $L$ de $A$ par des objets de $\mathcal{C}_1$, on ait $Z_d(L) \in \mathcal{C}_1$
02554: ($Z_d$ designe les cycles de dimension $d$).
02555: 
02556: Dans ces conditions, l'homomorphisme canonique de $K(\mathcal{C}'_1)$ dans $K(\mathcal{C}_1)$
02557: est un isomorphisme.
02558: \end{theorem}
02559: 
02560: \textit{Principe de la demonstration:} soit $A \in \mathcal{C}_1$; il existe en vertu de
02561: (ii) et (iii) une resolution finie $L$ de $A$ par des objets de $\mathcal{C}'_1$; posons
02562: 
02563: 
02564: \clearpage
02565: % ===== Source block: SGA 6 pages 051 100 =====
02566: \phantomsection
02567: \addcontentsline{toc}{section}{SGA 6 pages 051 100}
02568: \section*{SGA 6 pages 051 100}
02569: \setcounter{page}{51}
02570: 
02571: \selectlanguage{french}
02572: 
02573: \noindent\textbf{Note :} Ce document couvre les pages 51 \`a 100 du PDF original de SGA 6, correspondant aux pages 44 \`a 93 de la pagination originale de l'Expos\'e 0 (A. Grothendieck) et du d\'ebut de l'Expos\'e I (L. Illusie).
02574: 
02575: \tableofcontents
02576: \newpage
02577: 
02578: %=======================================================================
02579: \section*{Expos\'e 0 (suite)}
02580: \addcontentsline{toc}{section}{Expos\'e 0 (suite)}
02581: %=======================================================================
02582: 
02583: \subsection*{Fin de la d\'emonstration du th\'eor\`eme 2.3}
02584: \addcontentsline{toc}{subsection}{Fin de la d\'emonstration du th\'eor\`eme 2.3}
02585: 
02586: Alors $f(A) = \sum_{i \geq 0} (-1)^i \gamma_{\cCoh_0}(L^i)$. On montre que l'\'el\'ement de $\uK(\cCoh_0)$ ainsi construit ne d\'epend pas de la r\'esolution $\underline{L}$ choisie : si on a une deuxi\`eme r\'esolution $\underline{L}'$, on ``coiffe'' $\underline{L}$ et $\underline{L}'$ par une troisi\`eme r\'esolution $\underline{L}''$, avec des homomorphismes \emph{surjectifs} de $\underline{L}''$ dans $\underline{L}$ et $\underline{L}'$ (ce qui est possible gr\^ace \`a (i) et (ii)) ; on est alors ramen\'e \`a prouver $E_{\cCoh_0}(\underline{L}) = E_{\cCoh_0}(\underline{L}'')$. Or la diff\'erence des deux est $E_{\cCoh_0}(\underline{N})$ o\`u $\underline{N}$ est le noyau de l'homomorphisme surjectif de $\underline{L}''$ sur $\underline{L}$, noyau qui est dans $\cCoh_0$ d'apr\`es (i). Or $\underline{N}$ est un complexe acyclique en toutes dimensions, d'o\`u $E_{\cCoh_0}(\underline{N}) = 0$, d'o\`u la conclusion (N.B. $E_{\cCoh_0}(\underline{L})$, pour un complexe fini \`a valeurs dans $\cCoh_0$, d\'esigne la caract\'eristique d'Euler-Poincar\'e $\sum_{i \geq 0} (-1)^i \gamma_{\cCoh_0}(L^i)$ de $\underline{L}$).
02587: 
02588: On prouve de fa\c{c}on analogue que la fonction $f : \cCoh_1 \to \uK(\cCoh_0)$ ainsi construite est \emph{additive}, donc d\'efinit un homomorphisme de $\uK(\cCoh_1)$ dans $\uK(\cCoh_0)$, et il est imm\'ediat que cet homomorphisme est inverse de l'homomorphisme canonique de $\uK(\cCoh_0)$ dans $\uK(\cCoh_1)$.
02589: 
02590: \hfill $cqfd.$
02591: 
02592: \begin{remarque}
02593: Pour v\'erifier (ii), il suffit de v\'erifier que $\cCoh_1$ est contenu dans une sous-classe $\cCoh_2$ de $\cCoh$ satisfaisant aux deux conditions :
02594: \begin{enumerate}
02595: \item[(ii a)] Tout $A \in \cCoh_2$ est isomorphe \`a un quotient d'un $L \in \cCoh_0$.
02596: \item[(ii b)] Si $u : A \to B$, avec $A, B \in \cCoh_2$, alors $\Ker u \in \cCoh_2$.
02597: \end{enumerate}
02598: \end{remarque}
02599: 
02600: \begin{corollaire}
02601: Soit $X$ un espace alg\'ebrique quasi-projectif. Soient $\uF(X)$ la cat\'egorie des faisceaux alg\'ebriques coh\'erents sur $X$, $\xi(X)$ la cat\'egorie des faisceaux alg\'ebriques coh\'erents localement libres (\'equivalente \`a la cat\'egorie des fibr\'es vectoriels sur $X$), $\uF_0(X)$ la cat\'egorie des faisceaux alg\'ebriques coh\'erents $F$ tels que pour tout $x \in X$, $F_x$ soit un $\cO_x$-module de dimension cohomologique finie (on a donc $\xi(X) \subset \uF_0(X) \subset \uF(X)$). Alors l'homomorphisme canonique
02602: \begin{equation}\label{eq:2.12}
02603: \uK(\xi(X)) \longrightarrow \uK(\uF_0(X))
02604: \end{equation}
02605: est un \emph{isomorphisme}.
02606: \end{corollaire}
02607: 
02608: On va v\'erifier les conditions du th\'eor\`eme 2.3. Les conditions (i) et (iii) r\'esultent de la d\'efinition de la dimension cohomologique et du fait qu'un module projectif de type fini sur l'anneau local noeth\'erien $\cO_x$ est libre. Pour v\'erifier (ii), on va d\'emontrer que $\cCoh = \uF(X)$ satisfait aux conditions (ii a) et (ii b) de la remarque. La deuxi\`eme est triviale ; il reste donc \`a prouver le r\'esultat suivant :
02609: 
02610: \begin{proposition}[Serre]\label{prop:2.4}
02611: Soit $X$ un espace alg\'ebrique quasi-projectif. Alors tout faisceau alg\'ebrique coh\'erent sur $X$ est isomorphe au quotient d'un faisceau alg\'ebrique localement libre.
02612: \end{proposition}
02613: 
02614: Cette proposition est vraie si $X$ est une partie ferm\'ee de l'espace projectif en vertu de [2], puisque, pour $n$ assez grand, $F(n)$ est engendr\'e par ses sections, ce qui signifie que $F$ est un quotient de $\cO(-n)^k$. Dans le cas g\'en\'eral, $X$ est une partie ouverte d'une partie ferm\'ee d'un espace projectif, et il suffit d'utiliser le r\'esultat suivant (d\^u \`a Cartier et Serre dans des cas particuliers) :
02615: 
02616: \begin{proposition}\label{prop:2.5}
02617: Soit $F$ un faisceau alg\'ebrique coh\'erent sur une partie ouverte $U$ d'un espace alg\'ebrique $X$. Alors $F$ est isomorphe \`a la restriction d'un faisceau alg\'ebrique coh\'erent sur $X$.
02618: \end{proposition}
02619: 
02620: En zornifiant sur les ouverts contenant $U$ sur lesquels on peut prolonger $F$, on est ramen\'e au cas o\`u $X$ est affine. Soit $\overline{F}$ le faisceau alg\'ebrique sur $X$ dont les sections sur un ouvert $V$ sont les sections de $F$ sur $U \cap V$. On constate facilement que $\overline{F}$ est quasi-coh\'erent (i.e. d\'efini sur tout ouvert affine $V$ --- donc ici sur $X$ tout entier --- par un module, de type fini ou non, sur l'anneau de coordonn\'ees de l'ouvert en question). C'est imm\'ediat si $U$ est lui-m\^eme affine, car $\overline{F}$ est alors d\'efini sur l'ouvert $V$ par $\Gamma(U \cap V, F)$ consid\'er\'e comme module sur $\Gamma(V, \cO_X)$. Dans le cas g\'en\'eral, $U$ est r\'eunion finie d'ouverts affines $U_i$, donc $\overline{F}$ est le noyau d'un homomorphisme \'evident de faisceaux quasi-coh\'erents :
02621: \[
02622: \sum_i \overline{F}|_{U_i} \longrightarrow \sum_{i,j} \overline{F}|_{U_i \cap U_j}
02623: \]
02624: donc est lui-m\^eme quasi-coh\'erent. Comme $X$ est suppos\'e affine, $\overline{F}$ est donc engendr\'e par ses sections. Comme $\overline{F}|_U = F$ est coh\'erent, il existe un nombre fini de sections de $\overline{F}$ qui engendrent $\overline{F}|_U$. Ces sections engendrent un sous-faisceau coh\'erent de $\overline{F}$ ; dont la restriction \`a $U$ est isomorphe \`a $F$.
02625: 
02626: \hfill $cqfd.$
02627: 
02628: Supposons maintenant $X$ quasi-projective non singuli\`ere. Alors, en vertu du th\'eor\`eme des syzygies, on a, avec les notations du corollaire au th\'eor\`eme 2.3, $\uF_0(X) = \uF(X)$. Dans ce cas, on a donc :
02629: \begin{equation}\label{eq:2.13}
02630: \uK(\xi(X)) = \uK(\uF(X)).
02631: \end{equation}
02632: 
02633: On peut exprimer ainsi cet isomorphisme :
02634: 
02635: \begin{corollaire}[au th\'eor\`eme 2.3]\label{cor:2.3.2}
02636: Soit $X$ une vari\'et\'e quasi-projective non-singuli\`ere. Pour tout groupe ab\'elien $A$, il y a correspondance biunivoque entre les fonctions additives $f$ de fibr\'es vectoriels sur $X$, \`a valeurs dans $A$, et les fonctions additives $g$ de faisceaux alg\'ebriques coh\'erents sur $X$, \`a valeurs dans $A$. \`A toute $f$ correspond l'unique $g$ telle que l'on ait :
02637: \begin{equation}\label{eq:2.14}
02638: f(E) = g(\cO_X(E))
02639: \end{equation}
02640: pour tout fibr\'e vectoriel $E$ sur $X$, en notant $\cO_X(E)$ le faisceau des germes de sections r\'eguli\`eres de $E$ sur $X$.
02641: \end{corollaire}
02642: 
02643: Pour calculer $g(F)$ pour un faisceau alg\'ebrique coh\'erent quelconque, on consid\`ere une r\'esolution finie de $F$ par des faisceaux de la forme $\cO_X(E_i)$, et l'on a alors :
02644: \begin{equation}\label{eq:2.14bis}
02645: g(F) = \sum_{i \geq 0} (-1)^i f(E_i).
02646: \end{equation}
02647: 
02648: En particulier, si $X$ est connexe, on a envisag\'e au \S\,1 une fonction additive $\widetilde{C}(E)$ du fibr\'e vectoriel $E$, \`a valeurs dans le groupe $\widetilde{A} = \widetilde{A}(X)$. On en d\'eduit une fonction additive de faisceaux coh\'erents, soit $\widetilde{C}(F)$, et en passant aux composantes, des fonctions $F \mapsto c^i(F) \in A^i(X)$. Les $c^i(F)$ sont appel\'ees les \emph{classes de Chern du faisceau coh\'erent} $F$ ; elles se calculent th\'eoriquement \`a partir des classes de Chern des fibr\'es vectoriels par (2.14~bis) (mais o\`u il faut passer au deuxi\`eme membre \`a l'\'ecriture multiplicative). Quant \`a la composante de $\widetilde{C}(F)$ suivant $\bZ$, c'est le \emph{rang} du faisceau $F$, qui correspond au rang d'un module sur un anneau int\`egre : si $X$ est affine (cas auquel on peut toujours se ramener), $F$ est d\'efini par un module de type fini sur l'anneau de coordonn\'ees de $X$, et l'on prend le rang de ce module \ldots
02649: 
02650: %=======================================================================
02651: \subsection*{\S\,3. G\'en\'eralit\'es fonctorielles sur $\uK(X)$}
02652: \addcontentsline{toc}{subsection}{\S\,3. G\'en\'eralit\'es fonctorielles sur K(X)}
02653: %=======================================================================
02654: 
02655: Si $X$ est un espace alg\'ebrique, on pose pour abr\'eger $\uK(X) = \uK(\uF(X))$, et lorsque $X$ est quasi-projectif non singulier, on identifie ce groupe \`a $\uK(\xi(X))$ en vertu du corollaire 2 au th\'eor\`eme 2.3. Comme $\uK(\xi(X))$ est non seulement un groupe ab\'elien, mais un $\lambda$-anneau, il s'ensuit par transport de structure que $\uK(X)$ est aussi un $\lambda$-anneau, donc est muni d'une structure multiplicative et d'applications $\lambda^i$. Il importe de d\'efinir ces structures directement, sans passer aux fibr\'es vectoriels, ce qui sera fait ci-dessous.
02656: 
02657: Si $E$ est un fibr\'e vectoriel sur $X$, on d\'esigne par $\gamma(E)$ sa classe dans $\uK(X)$ ; si $F$ est un faisceau alg\'ebrique coh\'erent sur $X$, $\gamma(F)$ d\'esigne de m\^eme sa classe dans $\uK(X)$. Si $Y$ est une sous-vari\'et\'e irr\'eductible de $X$, on pose $\gamma(Y) = \gamma(\cO_Y)$, o\`u $\cO_Y$ d\'esigne le faisceau des anneaux locaux de $Y$, consid\'er\'e comme faisceau alg\'ebrique coh\'erent sur $X$. On d\'efinit par lin\'earit\'e le symbole $\gamma(Y)$ lorsque $Y$ est un cycle quelconque sur $X$. Quand il y aura des risques de confusion, on \'ecrira $\gamma_X$ pour $\gamma$.
02658: 
02659: Soit $X$ un espace alg\'ebrique quelconque. On introduit sur $\uK(X)$ une \emph{filtration}, en d\'esignant par $\uK(X)^i$ le sous-groupe de $\uK(X)$ engendr\'e par les $\gamma(F)$ avec $\dim \supp F \leq n - i$ (o\`u $n$ est la dimension de $X$). Il r\'esulte du ``lemme de d\'evissage'' de Serre, sous la forme donn\'ee dans [3], le r\'esultat suivant :
02660: 
02661: \begin{proposition}\label{prop:2.6}
02662: Le groupe $\uK(X)^i$ est engendr\'e par les $\gamma(Y)$, o\`u $Y$ parcourt l'ensemble des sous-vari\'et\'es de dimension $\leq n - i$ de $X$. Si $X$ est irr\'eductible, on a $G^0(\uK(X)) = \bZ$.
02663: \end{proposition}
02664: 
02665: Ce dernier isomorphisme est donn\'e par l'homomorphisme de $\uK(X)$ dans $\bZ$ d\'efini \`a l'aide de la notion de \emph{rang} d'un faisceau alg\'ebrique coh\'erent.
02666: 
02667: Supposons maintenant $X$ non singuli\`ere. S'inspirant des d\'eveloppements du d\'ebut du \S\,2, on d\'efinit une multiplication dans $\uK(X)$ en posant
02668: \begin{equation}\label{eq:2.15}
02669: \gamma(F) \cdot \gamma(G) = \sum_{i \geq 0} (-1)^i \gamma(\Tor_i(F,G))
02670: \end{equation}
02671: (les $\Tor_i$ \'etant pris, bien entendu, par rapport aux anneaux locaux de $X$). La suite exacte des $\Tor$ montre que la formule (2.15) d\'efinit bien une application biadditive de $\uK(X) \times \uK(X)$ dans $\uK(X)$ ; l'associativit\'e r\'esulte de la suite spectrale des $\Tor$ multiples. On v\'erifie aussi de cette fa\c{c}on que l'on peut calculer le produit de plusieurs facteurs par la formule
02672: \begin{equation}\label{eq:2.16}
02673: \gamma(F_1) \cdots \gamma(F_r) = \sum_{i \geq 0} (-1)^i \gamma(\Tor_i(F_1, \ldots, F_r))
02674: \end{equation}
02675: o\`u, dans le deuxi\`eme membre, les $\Tor_i$ sont des ``$\Tor$ simultan\'es'', calcul\'es par r\'esolution projective simultan\'ee des arguments $F_i$.
02676: 
02677: Si l'on suppose dans (2.15) que $G$ est localement libre, on trouve
02678: \begin{equation}\label{eq:2.15bis}
02679: \gamma(F) \cdot \gamma(G) = \gamma(F \otimes G),
02680: \end{equation}
02681: ce qui montre en particulier que $\gamma(\cO_X)$ est \'el\'ement unit\'e de $\uK(X)$, et que l'homomorphisme canonique de $\uK(\xi(X))$ dans $\uK(X)$ est un homomorphisme d'anneaux. La formule (2.15) est \'evidemment sugg\'er\'ee par la ``formule des $\Tor$'' de Serre en th\'eorie des intersections. Cette formule prouve la deuxi\`eme partie de la proposition suivante, la premi\`ere r\'esultant imm\'ediatement de la formule (2.15) et d'un calcul de $\Tor$ locaux :
02682: 
02683: \begin{proposition}\label{prop:2.7}
02684: Soient $X$ une vari\'et\'e non singuli\`ere, $Y$ et $Z$ deux cycles dans $X$. Si $Y$ et $Z$ se coupent partout transversalement, on a
02685: \begin{equation}\label{eq:2.17}
02686: \gamma(Y) \cdot \gamma(Z) = \gamma(Y \cdot Z)
02687: \end{equation}
02688: Si $Y$ et $Z$ sont homog\`enes et de dimension respectivement $n - i$ et $n - j$, et s'ils se coupent sans composante exc\'edentaire, on a
02689: \begin{equation}\label{eq:2.18}
02690: \gamma(Y) \cdot \gamma(Z) \equiv \gamma(Y \cdot Z) \pmod{\uK(X)^{i+j+1}}.
02691: \end{equation}
02692: \end{proposition}
02693: 
02694: On fera attention \`a ce qu'on n'a pas en g\'en\'eral \'egalit\'e dans (2.18).
02695: 
02696: On verra au \S\,4 que, si $X$ est quasi-projective non singuli\`ere, la structure d'anneau de $\uK(X)$ est compatible avec sa filtration et que l'anneau gradu\'e associ\'e $G(\uK(X))$ s'identifie \`a un quotient de $A(X)$ par un sous-groupe de torsion (peut-\^etre toujours r\'eduit \`a $0$~(*)).
02697: 
02698: Je ne sais d\'efinir directement les $\lambda^p(F)$ en termes de faisceaux que si le corps $k$ est de caract\'eristique $0$ (c'est la raison qui nous a emp\^ech\'e de d\'emontrer le th\'eor\`eme de Riemann-Roch lorsque $k$ n'est pas de caract\'eristique $\neq 0$). Soit $F$ un faisceau alg\'ebrique coh\'erent sur $X$ sur lequel op\`ere le groupe $\mathfrak{S}_p$ des permutations de $p$ lettres ; on note $\underline{\mathfrak{S}}_p(F)$ le faisceau des invariants de $F$ sous $\mathfrak{S}_p$.
02699: 
02700: % [Pages 58--60: Technical content on lambda operations, functoriality]
02701: 
02702: \begin{proposition}\label{prop:2.8}
02703: Soit $Y$ une sous-vari\'et\'e non singuli\`ere de $X$, $i$ le morphisme d'injection. Alors $i_! : \uK(Y) \to \uK(X)$ est donn\'e par la formule
02704: \begin{equation}\label{eq:2.19}
02705: i_!(\gamma_Y(F)) = \gamma_X(i_*(F))
02706: \end{equation}
02707: o\`u, dans le deuxi\`eme membre, $F$ est consid\'er\'e comme faisceau sur $Y$ (et donc $i_*(F)$ comme faisceau sur $X$, nul en dehors de $Y$).
02708: \end{proposition}
02709: 
02710: \begin{proposition}\label{prop:2.9}
02711: Soient $X$ une vari\'et\'e quasi-projective, $Y$ une partie ferm\'ee de $X$, $U = X \setminus Y$ son compl\'ementaire. Alors la suite
02712: \begin{equation}\label{eq:2.30}
02713: \uK(Y) \longrightarrow \uK(X) \longrightarrow \uK(U) \longrightarrow 0
02714: \end{equation}
02715: est exacte.
02716: \end{proposition}
02717: 
02718: N.B. On comparera avec la proposition 2.3.
02719: 
02720: Le fait que $\uK(X) \to \uK(U)$ soit surjectif r\'esulte aussit\^ot de la proposition 2.5 (ou aussi de la proposition 2.6). Reste \`a prouver que si $x \in \uK(X)$ a une restriction nulle \`a $\uK(U)$, il est dans l'image de $\uK(Y)$. Pour ceci, on revient \`a la d\'efinition explicite de $\uK(X)$ et de $\uK(U)$ : on a deux faisceaux alg\'ebriques coh\'erents $F$ et $G$ sur $X$ avec $x = \gamma_X(F) - \gamma_X(G)$ et $\gamma_U(F) - \gamma_U(G) = 0$ ; cette derni\`ere relation signifie qu'on peut trouver des faisceaux coh\'erents $P'$ et $Q'$ sur $U$, munis de filtrations dont les facteurs soient deux \`a deux isomorphes sur $U$, et tels que l'on ait un isomorphisme sur $U$ :
02721: \[
02722: P = F + Q' \xrightarrow{\quad\sim\quad} Q = G + P'.
02723: \]
02724: En vertu de la proposition 2.5, les faisceaux $P'$ et $Q'$ sont des restrictions \`a $U$ de faisceaux coh\'erents sur $X$ (not\'es par les m\^emes lettres) ; de m\^eme, leurs filtrations proviennent de filtrations sur les faisceaux correspondants sur $X$, en vertu du lemme 2.11 ci-dessous.
02725: 
02726: \begin{lemme}\label{lem:2.11}
02727: Soient $X$ un espace alg\'ebrique, $F$ un faisceau alg\'ebrique coh\'erent sur $X$, $U$ une partie ouverte de $X$ et $G$ un sous-faisceau coh\'erent de $F|_U$. Alors $G$ est la restriction \`a $U$ d'un sous-faisceau coh\'erent de $F$.
02728: \end{lemme}
02729: 
02730: Soit $\widetilde{G}$ le sous-faisceau de $F$ dont les sections sur l'ouvert $V$ de $X$ sont les sections de $F$ sur $V$ dont la restriction \`a $U \cap V$ est une section de $G$. Tout revient \`a prouver que ce faisceau est coh\'erent. Pour ceci, on peut supposer $X$ et $U$ affines (car $U$ est r\'eunion d'ouverts affines $U_i$, et $G$ est alors intersection des faisceaux analogues $\widetilde{G}|_{U_i}$). Alors l'anneau de coordonn\'ees de $U$ est \ldots
02731: 
02732: %=======================================================================
02733: \subsection*{\S\,4. Quelques r\'esultats techniques}
02734: \addcontentsline{toc}{subsection}{\S\,4. Quelques r\'esultats techniques}
02735: %=======================================================================
02736: 
02737: Soit $X$ un espace alg\'ebrique ; on a vu au \S\,3 que la projection $p$ de $X \times \bP_k^r$ sur $X$ d\'efinit un homomorphisme $p^!$ de $\uK(X)$ dans $\uK(X \times \bP_k^r)$. Nous allons prouver que :
02738: 
02739: \begin{proposition}\label{prop:2.10}
02740: $p^!$ est surjectif.
02741: \end{proposition}
02742: 
02743: On proc\'ede par r\'ecurrence sur $r$. Pour $r = 0$, c'est trivial. Supposons l'assertion v\'erifi\'ee pour $r - 1$, et prouvons-la pour $r$. Il suffit de v\'erifier que pour toute partie ferm\'ee irr\'eductible $Y$ de $X \times \bP_k^r$, $\gamma(Y)$ est dans l'image de $p^!$. Il suffit de continuer le raisonnement en supposant $X$ irr\'eductible ; si $p(Y)$ n'est pas dense dans $X$, notre assertion r\'esulte de l'hypoth\`ese de r\'ecurrence ; sinon, on a $Y = X \times \bP_k^r$ et l'assertion est triviale, ou bien $Y$ est une hypersurface dans $X \times \bP_k^r$. Soit alors $U$ un ouvert affine de $X$ form\'e de points non singuliers ; si $D$ est le diviseur sur $U \times \bP_k^r$ induit par $Y$, on sait (alg\`ebre commutative) qu'il est lin\'eairement \'equivalent \`a un diviseur de la forme $D' \times \bP_k^r$, o\`u $D'$ est un diviseur sur $U$. Alors $\gamma_{U \times \bP_k^r}(D) - \gamma_{U \times \bP_k^r}(D' \times \bP_k^r)$ appartient \`a $\uK(U)^1$ (corollaire de la proposition 2.10), et d'apr\`es ce qu'on a d\'ej\`a d\'emontr\'e, cet \'el\'ement est de la forme $p^*(x')$ avec $x' \in \uK(U)$. Donc $\gamma_{U \times \bP_k^r}(D) = p^*(y')$ avec $y' = x' + \gamma_U(D')$ ; mais $y'$ est restriction d'un \'el\'ement $y$ de $\uK(X)$ (prop. 2.10), et par suite la restriction de $\gamma_{X \times \bP_k^r}(Y) - p^!(y)$ \`a $U \times \bP_k^r$ est nulle ; d'apr\`es la proposition 2.10, $\gamma_{X \times \bP_k^r}(Y) - p^!(y)$ provient de $\uK(Z)$, o\`u $Z$ est une partie ferm\'ee de $X \times \bP_k^r$ ne rencontrant pas $U \times \bP_k^r$. La projection de $Z$ sur $X$ est donc une partie ferm\'ee ne rencontrant pas $U$, donc contenue dans $X \setminus U$, et on conclut gr\^ace \`a l'hypoth\`ese de r\'ecurrence sur $r$.
02744: 
02745: \begin{corollaire}\label{cor:2.10.1}
02746: Supposons $X$ connexe de dimension $n$, sans singularit\'es en dimension $n - 1$. Si $D$ et $D'$ sont deux diviseurs lin\'eairement \'equivalents sur $X$, on a $\gamma(D) \equiv \gamma(D') \pmod{\uK(X)^2}$.
02747: \end{corollaire}
02748: 
02749: En vertu du corollaire 1, on peut supposer $X$ non singuli\`ere ; le corollaire r\'esulte alors de la formule g\'en\'erale :
02750: \begin{equation}\label{eq:2.31}
02751: \gamma(D) \equiv 1 - \gamma(\cL(-D)) \pmod{\uK(X)^2},
02752: \end{equation}
02753: o\`u $\cL(-D)$ d\'esigne le fibr\'e vectoriel de rang 1 d\'efini par le diviseur $-D$. Pour v\'erifier (2.31), on l'\'ecrit sous la forme
02754: \begin{equation}\label{eq:2.31bis}
02755: \gamma(\cL(-D)) \equiv 1 - \gamma(D) \pmod{\uK(X)^2},
02756: \end{equation}
02757: et l'on note que les deux membres sont des homomorphismes du groupe des diviseurs dans le groupe multiplicatif des \'el\'ements inversibles de l'anneau $\uK(X)/\uK(X)^2$ (noter que $\uK(X)/\uK(X)^2$ est de carr\'e nul, ce qui se v\'erifie tr\`es facilement sur la formule explicite (2.15), en utilisant le corollaire 1). On est ainsi ramen\'e au cas o\`u $D$ est une hypersurface irr\'eductible. Plus g\'en\'eralement, si $D$ est une hypersurface dont les composantes sont disjointes, on a l'\'egalit\'e :
02758: \begin{equation}\label{eq:2.32}
02759: \gamma(D) = 1 - \gamma(\cL(-D)).
02760: \end{equation}
02761: 
02762: \begin{theoreme}\label{thm:2.12}
02763: Soit $X$ une vari\'et\'e quasi-projective non singuli\`ere. Soient $x \in \uK(X)$, $y \in \uK(X)$ et deux points $a, a' \in k^*$ tels que l'on ait :
02764: \[
02765: z = f_a^!(y), \quad z = f_{a'}^!(y),
02766: \]
02767: o\`u l'on pose $f_a^!(x) = (x, a)$. En vertu de la formule (2.23), on a alors
02768: \[
02769: \gamma_Y(z) \equiv f^!(\gamma_X(y)) \pmod{\uK(X)^{p+1}}.
02770: \]
02771: \end{theoreme}
02772: 
02773: %=======================================================================
02774: \subsection*{\S\,5. D\'efinition faisceautique des classes de Chern. Application \`a l'\'etude des morphismes d'injection}
02775: \addcontentsline{toc}{subsection}{\S\,5. D\'efinition faisceautique des classes de Chern}
02776: %=======================================================================
02777: 
02778: \begin{theoreme}\label{thm:2.16}
02779: Soit $X$ une vari\'et\'e quasi-projective non singuli\`ere. Alors, pour tout $x \in \uK(X)$, et pour tout entier $i \geq 0$, l'\'el\'ement $\gamma^i(x - \varepsilon(x))$ est de filtration $\geq i$. Soit $c^i(x) \in G^i(\uK(X))$ la classe de cet \'el\'ement modulo $\uK(X)^{i+1}$ ; si l'on pose :
02780: \[
02781: \widetilde{c}(x) = [\varepsilon(x), 1 + \sum_{i \geq 0} c^i(x)] \in \widetilde{G}(\uK(X)),
02782: \]
02783: alors $\widetilde{c}$ satisfait aux conditions analogues \`a celles du th\'eor\`eme 2.2.
02784: \end{theoreme}
02785: 
02786: La premi\`ere assertion r\'esulte du th\'eor\`eme 2.14, ainsi que le fait que $x \mapsto \widetilde{c}(x)$ soit un homomorphisme de $\lambda$-anneaux (la formule (1.25) implique $\gamma_t(x+y) = \gamma_t(x)\,\gamma_t(y)$, ce qui montre d\'ej\`a que $\widetilde{c}$ est un homomorphisme additif). La ``fonctorialit\'e'' de $\widetilde{c}$ se v\'erifie imm\'ediatement, et la ``compatibilit\'e avec les structures multiplicatives'' se prouve \`a l'aide de la formule (2.15) (qui donne l'expression explicite du produit dans $\uK(X)$) et d'un calcul de $\Tor$ locaux (qui permet de comparer les deux membres de la condition (ii bis) du th\'eor\`eme 2.2).
02787: 
02788: \hfill $cqfd.$
02789: 
02790: \begin{corollaire}\label{cor:2.16.1}
02791: La filtration de $\uK(X)$ est compatible avec sa structure d'anneau, i.e. on a $\uK(X)^i \cdot \uK(X)^j \subset \uK(X)^{i+j}$ quels que soient les entiers $i, j \geq 0$. On a de plus un homomorphisme surjectif naturel d'anneaux gradu\'es
02792: \begin{equation}\label{eq:2.34}
02793: \alpha : A(X) \longrightarrow G(\uK(X))
02794: \end{equation}
02795: o\`u le deuxi\`eme membre d\'esigne l'anneau gradu\'e associ\'e \`a $\uK(X)$ muni de la filtration ci-dessus, d\'efini par la condition
02796: \[
02797: \alpha(Z^p) = \gamma(Z^p) \pmod{\uK(X)^{p+1}}.
02798: \]
02799: \end{corollaire}
02800: 
02801: Ce corollaire r\'esulte imm\'ediatement du th\'eor\`eme 2.12, des propositions 2.6 et 2.7 et du th\'eor\`eme de Chow 2.1, par r\'ecurrence sur l'entier $i + j$.
02802: 
02803: \begin{corollaire}\label{cor:2.16.2}
02804: Soient $X$ et $Y$ deux vari\'et\'es quasi-projectives non singuli\`eres et $f$ un morphisme de $X$ dans $Y$. Alors $f^! : \uK(Y) \to \uK(X)$ est compatible avec les filtrations de $\uK(Y)$ et $\uK(X)$, et l'homomorphisme gradu\'e correspondant $G(f^!) : G(\uK(Y)) \to G(\uK(X))$ rend commutatif le diagramme
02805: \[
02806: \begin{array}{ccc}
02807: A(Y) & \xrightarrow{\quad f^*\quad} & A(X) \\
02808: \downarrow & & \downarrow \\
02809: G(\uK(Y)) & \xrightarrow{\quad G(f^!)\quad} & G(\uK(X))
02810: \end{array}
02811: \]
02812: \end{corollaire}
02813: 
02814: Utilisant le graphe de $f$, on d\'ecompose $f$ en produit d'une immersion de $X$ dans $X \times Y$ et d'une projection de $X \times Y$ sur $Y$. Pour cette derni\`ere le corollaire est trivial, et pour une immersion, on proc\`ede comme pour le corollaire 1, en utilisant la Proposition 2.8 \`a la place de la proposition 2.7.
02815: 
02816: \begin{proposition}[``d\'ecomposition cellulaire'']\label{prop:2.13}
02817: Soit $P$ un espace alg\'ebrique (irr\'eductible ou non), muni d'une famille croissante $P_0 \subset P_1 \subset \cdots \subset P_n = P$ de parties ferm\'ees, telles que pour tout $i$ avec $0 \leq i < n$, les composantes connexes de $P^i - P^{i-1}$ soient des espaces alg\'ebriques isomorphes \`a $k^{n_i}$ (on pose $P^{-1} = \emptyset$). Alors :
02818: \begin{enumerate}
02819: \item[a)] Pour tout espace alg\'ebrique $X$, l'application naturelle de $\uK(X) \otimes \uK(P)$ dans $\uK(X \times P)$ d\'eduite du produit tensoriel de faisceaux sur $X$ et $P$, est surjective.
02820: \item[b)] Les $\gamma_{P^i}$ forment un syst\`eme de g\'en\'erateurs du groupe $\uK(P)$.
02821: \end{enumerate}
02822: \end{proposition}
02823: 
02824: La d\'emonstration se fait par r\'ecurrence sur $n$, en utilisant le corollaire de la proposition 2.9 et la proposition 2.10. (N.B. Malheureusement, j'ignore si l'application naturelle de $\uK(X) \otimes \uK(P)$ dans $\uK(X \times P)$ est toujours bijective.)
02825: 
02826: %=======================================================================
02827: \subsection*{Th\'eor\`eme de Riemann-Roch}
02828: \addcontentsline{toc}{subsection}{Th\'eor\`eme de Riemann-Roch}
02829: %=======================================================================
02830: 
02831: Soit $f : Y \to X$ un morphisme propre d'une vari\'et\'e quasi-projective non singuli\`ere $Y$ dans une autre $X$, et soit $y \in \uK(Y)$. On d\'efinit :
02832: \begin{equation}\label{eq:rr}
02833: f_!(\ch(y)\,T(Y)) = \ch(f_!(y))\,T(X) \quad \text{dans } A(X) \otimes \bQ.
02834: \end{equation}
02835: 
02836: \textbf{Th\'eor\`eme de Riemann-Roch.} Soit $f : Y \to X$ un morphisme propre d'une vari\'et\'e quasi-projective non singuli\`ere $Y$ dans une autre $X$, et soit $y \in \uK(Y)$. On a :
02837: \begin{equation}\label{eq:rr.1}
02838: f_*(\ch(y)\,T(Y)) = \ch(f_!(y))\,T(X) \quad \text{dans } A(X) \otimes \bQ.
02839: \end{equation}
02840: 
02841: \textbf{D\'emonstration} (abr\'eg\'ee).
02842: 
02843: \begin{lemme}
02844: Consid\'erons des morphismes propres de vari\'et\'es $Z \xrightarrow{g} Y \xrightarrow{f} X$. Si le th\'eor\`eme R.R. est vrai pour $f$ et $g$, il l'est pour $f \circ g$. S'il est vrai pour $fg$ et $g$, il l'est pour $f$. Si $f_*$ est injectif, et si le th\'eor\`eme R.R. est vrai pour $fg$ et $f$, il l'est pour $g$.
02845: \end{lemme}
02846: 
02847: Trivial. (En fait, lorsqu'on suppose R.R. vrai pour $fg$ et $g$, il faut encore supposer que $g_*$ est surjectif pour pouvoir affirmer que R.R. est vrai pour $f$.)
02848: 
02849: Un morphisme propre \'etant compos\'e d'une injection et d'une projection $X \times \bP^r \to X$ (o\`u $\bP^r$ est un espace projectif), on est ramen\'e \`a traiter s\'epar\'ement ces deux cas. Le deuxi\`eme, on l'a vu, se ram\`ene \`a du calcul standard, moyennant la proposition 2.13 (``d\'ecomposition cellulaire'') de R.R.R. On va donc supposer $Y \subset X$, avec le morphisme d'injection $i$. La formule (1) devient alors :
02850: \begin{equation}\label{eq:rr.2}
02851: \ch(i_!(y)) = i_*(\ch(y)\,T(E)^{-1}),
02852: \end{equation}
02853: o\`u $E = T_{X/Y}$ est le fibr\'e normal \`a $Y$ dans $X$.
02854: 
02855: \begin{lemme}
02856: Supposons qu'il existe sur $X$ un fibr\'e vectoriel de rang $p = \codim_X Y$, tel que
02857: \[
02858: \lambda_{-1}(E^*) = \sum (-1)^p \lambda^p(E^*) \text{ et } \lambda_{-1}(E^*) = 0.
02859: \]
02860: Si alors on a $y = i^*(x)$, avec $x \in \uK(X)$, la formule (2) est valable.
02861: \end{lemme}
02862: 
02863: D\'emonstration imm\'ediate, \`a partir du fait que $\ch$ est un homomorphisme d'anneaux, et que l'on a
02864: \begin{equation}\label{eq:rr.5}
02865: \ch(\lambda_{-1}(E^*)) = c^p(E)\,T(E)^{-1}.
02866: \end{equation}
02867: 
02868: \begin{corollaire}
02869: R.R. est vrai si $Y$ est une hypersurface et si $y$ est de la forme $i^*(x)$, avec $x \in \uK(X)$.
02870: \end{corollaire}
02871: 
02872: \begin{lemme}
02873: Supposons que $X = Y \times P$, $P$ \'etant un espace projectif, et que $i$ soit l'injection $u \mapsto (u, v_0)$, $v_0 \in P$. Alors $i_*$ est injectif et R.R. est vrai.
02874: \end{lemme}
02875: 
02876: C'est du pur calcul standard.
02877: 
02878: Supposons de nouveau $X$, $Y$, $y$ quelconques (toujours avec $Y \subset X$). Soit $X'$ la vari\'et\'e \'eclat\'ee de $X$ \`a l'aide de $Y$, et $Y'$ l'hypersurface image r\'eciproque de $Y$ :
02879: 
02880: On a les formules (o\`u $p = \codim_X Y$) :
02881: \begin{align}
02882: i'_!(y) &= j_!(y \cdot \lambda_{-1}(F^*)), \tag{6} \\
02883: \lambda_{-1}(F^*) &\equiv \xi^{p-1}(E) \pmod{(1 - \cL)}, \tag{7} \\
02884: f_!(1) &= 1, \quad f_! f^! = 1, \tag{8} \\
02885: g^* i_!(a) &= j_!(a \cdot c^p(E)) \pmod{\Ker f_!} \quad \text{pour tout } a \in A(Y). \tag{9}
02886: \end{align}
02887: (Pour simplifier l'\'ecriture, on a fait op\'erer $\uK(Y)$ sur $\uK(Y')$ et $A(Y)$ sur $A(Y')$ gr\^ace aux homomorphismes $g^!$ et $g^*$. Dans la formule (7), il semble que $\xi^{p-1}(E)$ d\'esigne $\lambda^{p-1}(E - 2)$, c'est-\`a-dire $\sum_{i=0}^{p-1} (-1)^i \lambda^{p-1-i}(E) \xi^i$. La congruence a lieu dans $\uK(Y')$.)
02888: 
02889: La formule (6) r\'esulte de l'\'egalit\'e
02890: \[
02891: \Tor_i^{\cO_X}(\cO_Y, \cO_{Y'}) = \lambda^i(F^*)
02892: \]
02893: et de la suite spectrale des $\Tor$.
02894: 
02895: La formule (7) r\'esulte des formules
02896: \[
02897: g_!(\cL^i) = 0 \text{ pour } 0 < i < p-1, \quad g_!(1) = 1
02898: \]
02899: elles-m\^emes cons\'equences de K\"unneth et de la cohomologie des espaces projectifs.
02900: 
02901: La formule (8) est triviale, enfin (9) r\'esulte de $g_!(c^{p-1}(F)) = c^p(E)$ et de la fonctorialit\'e de $c^p$.
02902: 
02903: De la formule (7) r\'esulte aussit\^ot $\gamma \cdot 1 \in G^0(\uK(Y'))$ pourvu que l'on ait
02904: \[
02905: \gamma \cdot \xi^{p-1}(E) \cdot c^p(E) \in i_!(A(Y)).
02906: \]
02907: Supposons cette condition (11) satisfaite. D'apr\`es le corollaire au lemme 2, R.R. est vrai pour $(Y', x', y \cdot \lambda_{-1}(F^*))$. D'autre part, pour prouver (2), il suffit, en vertu de la formule (8), de prouver que
02908: \[
02909: f^*\ch(i_!(y)) = f^*\ch(i_!(y)\,T(E)^{-1}) \pmod{\Ker f_!}.
02910: \]
02911: Or, cette derni\`ere formule se v\'erifie aussit\^ot, utilisant (6), (9), R.R. pour $(Y', X', y \cdot \lambda_{-1}(F^*))$ et la formule :
02912: \begin{equation}\label{eq:rr.13}
02913: \ch(\lambda_{-1}(F^*)) = c^p(F)\,T(F)^{-1} \quad \text{(cf. formule (5)).}
02914: \end{equation}
02915: Ainsi, R.R. est prouv\'e chaque fois que $\gamma \cdot c^p(E) = 0$, en particulier (puisque $\gamma \cdot c^p(E)$ est de filtration $\geq p-1$ dans $\uK(Y)$ filtr\'e par la codimension des supports) chaque fois que $p-1 > \dim Y$, i.e. chaque fois que $\dim Y < \dim X - 1$.
02916: 
02917: Or, on se ram\`ene \`a ce cas, en conjuguant \`a un plongement projectif convenable. De fa\c{c}on pr\'ecise, soient $P$ un espace projectif, $v_0$ un point de $P$, $j$ l'injection $y \mapsto (y, v_0)$ de $Y$ dans $Y \times P$. Soit $Y'$ une sous-vari\'et\'e non singuli\`ere de $X \times P$ de m\^eme dimension que $Y$, contenant $j(Y)$, et telle que $Y' \cap (X \times \{v_0\}) = j(Y)$ transversalement. Alors $j_!(y) = j_!(y')$, o\`u $y'$ est un \'el\'ement de $\uK(Y')$ dont l'image directe (pour l'injection $Y' \to X \times P$) est l'image directe de $y$ (pour $Y \to X$, compos\'ee avec $X \to X \times \{v_0\}$). On est ainsi ramen\'e \`a prouver R.R. pour l'injection de $Y'$ dans $X \times P$, ce qui nous ram\`ene au cas o\`u $\dim Y' < \dim(X \times P) - 1$.
02918: 
02919: Dans ce deuxi\`eme cas, je n'ai pas fait la d\'etermination compl\`ete, il manque un petit quelque chose. Avec les notations du petit papier, il faut prouver la formule (9) (mais comme \'egalit\'e) :
02920: \begin{equation}\label{eq:rr.9}
02921: g^*(i_!(a)) = j_!(a \cdot c^p(E)) \quad a \in A(Y).
02922: \end{equation}
02923: On constate que les deux membres ont m\^eme image par $f_!$, et m\^eme image par $j^*$, de sorte que la question est \'equivalente \`a
02924: \begin{equation}\label{eq:rr.9bis}
02925: \Ker f_! \cap \Ker j^* = 0.
02926: \end{equation}
02927: En tout cas, (9) est vraie en r\'eduisant \`a $G(\uK(X'))$ (ce qui d\'etruit un peu de torsion, tout au plus), car elle r\'esulte alors de (6) ; il n'y a donc gu\`ere de doute qu'elle ne soit vraie.
02928: 
02929: Je te signale d'ailleurs que l'\'egalit\'e (9) r\'esulte tr\`es facilement de (6) sans passer par l'ennuyeuse (6~bis) \ldots
02930: 
02931: \newpage
02932: %=======================================================================
02933: % PART II: Exposé I by L. Illusie
02934: %=======================================================================
02935: \section*{Expos\'e I --- G\'en\'eralit\'es sur les conditions de finitude dans les cat\'egories d\'eriv\'ees}
02936: \addcontentsline{toc}{section}{Expos\'e I --- Conditions de finitude dans les cat\'egories d\'eriv\'ees (L. Illusie)}
02937: \begin{center}\textbf{par L. ILLUSIE}\end{center}
02938: 
02939: %=======================================================================
02940: \subsection*{0. Introduction}
02941: \addcontentsline{toc}{subsection}{0. Introduction}
02942: %=======================================================================
02943: 
02944: L'objet des Expos\'es I \`a IV est de d\'evelopper, avec la g\'en\'eralit\'e convenant dans ce S\'eminaire, le formalisme des conditions de finitude dans les cat\'egories d\'eriv\'ees des topos annel\'es.
02945: 
02946: Comme il a \'et\'e dit dans Exp.~0, c'est le besoin de d\'efinir sur des sch\'emas arbitraires des ``groupes de Grothendieck'' poss\'edant de bonnes propri\'et\'es de variance qui a oblig\'e \`a g\'en\'eraliser et \`a assouplir les notions de finitude utilis\'ees jusqu'\`a pr\'esent.
02947: 
02948: Une notion classique comme celle de faisceau coh\'erent sur un espace annel\'e $(X, \cO_X)$ devient sans int\'er\^et d\`es que $\cO_X$ n'est plus coh\'erent. On pourrait songer \`a la remplacer par la notion de pr\'esentation finie, mais celle-ci pr\'esente l'inconv\'enient que le noyau d'un \'epimorphisme de modules de pr\'esentation finie n'est pas en g\'en\'eral de pr\'esentation finie.
02949: 
02950: Si $X$ est un espace annel\'e en anneaux locaux noeth\'eriens, un $\cO_X$-module coh\'erent $F$ admet localement une r\'esolution finie gauche par des modules libres de type fini, i.e. il existe localement une suite exacte
02951: \[
02952: 0 \longrightarrow L_n \longrightarrow L_{n-1} \longrightarrow \cdots \longrightarrow L_0 \longrightarrow F \longrightarrow 0
02953: \]
02954: o\`u les $L_i$ sont libres de type fini. Le faisceau $\cO_X$ n'\'etant plus suppos\'e coh\'erent, disons qu'un $\cO_X$-module $F$ est \emph{pseudo-coh\'erent} s'il v\'erifie la condition pr\'ec\'edente, i.e. est de $n$-pr\'esentation finie pour tout $n$. Alors, la notion de pseudo-coh\'erence poss\`ede, vis-\`a-vis des suites exactes courtes, la m\^eme propri\'et\'e de stabilit\'e que la notion de coh\'erence : si deux des termes d'une suite exacte courte de $\cO_X$-modules sont pseudo-coh\'erents, le troisi\`eme l'est aussi. On g\'en\'eralise ais\'ement la notion pr\'ec\'edente aux complexes de faisceaux. Un complexe $F$ de $\cO_X$-modules est dit \emph{pseudo-coh\'erent} si chaque $F^n$ est pseudo-coh\'erent en tant que complexe ou en tant que module.
02955: 
02956: La notion de complexe pseudo-coh\'erent a d'excellentes vertus de stabilit\'e, d\'ecrites dans Exp.~I \S\,1 et 2. D'abord c'est une notion locale. Ensuite, si $f : (X, \cO_X) \to (Y, \cO_Y)$ est un morphisme d'espaces annel\'es, l'image inverse $Lf^*$ d'un complexe pseudo-coh\'erent sur $Y$ est pseudo-coh\'erente sur $X$, pourvu que le tore dimension finie rel. \`a $f$ (i.e. rel. au faisceau d'anneaux $f^{-1}(\cO_Y)$).
02957: 
02958: Un complexe coh\'erent $F$ admet localement une r\'esolution finie gauche par des modules libres de type fini, i.e. il existe localement une suite exacte
02959: \[
02960: 0 \longrightarrow L_n \longrightarrow L_{n-1} \longrightarrow \cdots \longrightarrow L_0 \longrightarrow F \longrightarrow 0
02961: \]
02962: o\`u les $L_i$ sont libres de type fini. Il est donc naturel, quand $\cO_X$ est un faisceau d'anneaux locaux quelconque, de s'int\'eresser aux modules qui jouissent de la propri\'et\'e pr\'ec\'edente. De tels modules sont appel\'es \emph{parfaits}. Plus g\'en\'eralement, on dit qu'un complexe $F$ de $\cO_X$-modules est \emph{parfait} s'il existe localement un quasi-isomorphisme $L \to F$, o\`u $L$ est un complexe \`a degr\'es born\'es et \`a composantes libres de type fini. On montre dans Exp.~I \S\,3 que la sous-cat\'egorie pleine $D(X)_{\text{parf}}$ de $D(X)$ form\'ee des complexes parfaits, est, de m\^eme que $D(X)_{\text{coh}}$, une sous-cat\'egorie triangul\'ee, ``stable'' par image inverse et produit tensoriel d\'eriv\'e.
02963: 
02964: Des r\'esolutions globales de complexes parfaits sont donn\'ees dans l'Exp.~II. On montre en particulier qu'il existe des r\'esolutions globales dans les cas suivants : $X$ est le topos zariskien d'un sch\'ema affine, ou plus g\'en\'eralement d'un sch\'ema quasi-compact poss\'edant un faisceau inversible ample, voire une famille ample de faisceaux inversibles au sens de Kleiman (par exemple un sch\'ema r\'egulier) ; ou encore $X$ est le topos des faisceaux d'ensembles sur un espace topologique compact, $\cO_X$ le faisceau des fonctions continues \`a valeurs complexes. \`A titre d'illustration --- et pour mettre en \'evidence la souplesse des notions introduites --- on donne en appendice une d\'efinition ``purement faisceautique'' de l'indice d'une famille d'op\'erateurs elliptiques.
02965: 
02966: L'Exp.~III \'etudie la stabilit\'e des conditions de finitude par image directe d\'eriv\'ee. Pour obtenir des r\'esultats, on est amen\'e \`a introduire la notion de complexe \emph{pseudo-coh\'erent relatif} \`a un morphisme $f : X \to Y$. Cela veut dire que $F$ est pseudo-coh\'erent rel. \`a $f$ et localement de tor-dimension finie rel. \`a $f$ (i.e. rel. au faisceau d'anneaux $f^{-1}(\cO_Y)$).
02967: 
02968: Le th\'eor\`eme central de l'Exp.~III est le \emph{th\'eor\`eme de finitude}, qui affirme (sous une forme un peu plus pr\'ecise) que, si $f : X \to Y$ est un morphisme propre de $S$-sch\'emas localement de type fini, le foncteur $Rf_*$ transforme complexes pseudo-coh\'erents rel. \`a $S$ en complexes pseudo-coh\'erents rel. \`a $S$ ; ce th\'eor\`eme est en r\'ealit\'e une conjecture, mais est prouv\'e n\'eanmoins dans les deux cas particuliers suivants :
02969: \begin{enumerate}
02970: \item[a)] $S$ est localement noeth\'erien ;
02971: \item[b)] $f$ est projectif ;
02972: \end{enumerate}
02973: il semble malheureusement que l'extension au cas g\'en\'eral soit du m\^eme ordre de difficult\'e que le th\'eor\`eme analogue en G\'eom\'etrie Analytique (``th\'eor\`eme'' de Grauert).
02974: 
02975: Combinant le th\'eor\`eme de finitude avec une formule essentiellement triviale dite \emph{formule de projection}, on obtient des crit\`eres maniables pour la stabilit\'e de la perfection relative par image directe. On retrouve, en corollaire, les ``th\'eor\`emes de continuit\'e et de semi-continuit\'e de Grauert'' (EGA III 7.6).
02976: 
02977: L'Exp.~IV traduit en langage de ``groupes de Grothendieck'' les r\'esultats des Exp.~I \`a III. Sur un topos annel\'e $(X, \cO_X)$, on note $K^*(X)$ (resp. $K_*(X)$) le ``groupe de Grothendieck'' de la cat\'egorie des complexes parfaits et de tor-dimension finie (resp. pseudo-coh\'erents et \`a cohomologie born\'ee).
02978: 
02979: 
02980: %=======================================================================
02981: \subsection*{1. Definitions preliminaires}
02982: \addcontentsline{toc}{subsection}{1. Definitions preliminaires}
02983: %=======================================================================
02984: 
02985: \subsubsection*{1.1. Categories fibrees a fibres additives (resp. abeliennes, resp. triangulees)}
02986: \addcontentsline{toc}{subsubsection}{1.1. Categories fibrees a fibres additives}
02987: 
02988: \begin{definition}
02989: Soit $S$ une categorie, et soit $\mathbf{C}$ une $S$-categorie fibree (SGA 1 VI 6.1) a fibres des categories additives (resp. abeliennes, resp. triangulees). Nous dirons que $\mathbf{C}$ est une $S$-categorie additive (resp. abelienne, resp. triangulee) (ou que $\mathbf{C}$ est additive (resp. abelienne, resp. triangulee) \emph{sur} $S$) si, pour toute fleche $f : X \to Y$ de $S$, le foncteur image inverse (determine a isomorphisme unique pres) $f^* : \mathbf{C}_Y \to \mathbf{C}_X$ est additif (resp. exact, resp. triangule).
02990: \end{definition}
02991: 
02992: \begin{definition}
02993: Soit $\mathbf{C}$ une $S$-categorie triangulee, et soit $\mathbf{C}'$ une sous-$S$-categorie pleine de $\mathbf{C}$. On dit que $\mathbf{C}'$ est une sous-$S$-categorie triangulee de $\mathbf{C}$ si, pour tout $X \in \operatorname{ob} S$, $\mathbf{C}'_X$ est une sous-categorie triangulee de $\mathbf{C}_X$.
02994: \end{definition}
02995: 
02996: \begin{definition}
02997: Soit $\mathbf{C}$ une $S$-categorie additive. On definit (SGA 4 XVII 2.2, 2.4) une $S$-categorie additive (resp. triangulee) $\mathbf{C}(\mathbf{C})$ (resp. $\mathbf{K}(\mathbf{C})$), dont la fibre en $X \in \operatorname{ob} S$ est la categorie $\mathbf{C}(\mathbf{C}_X)$ (resp. $\mathbf{K}(\mathbf{C}_X)$) des complexes (resp. complexes ``a homotopie pres'') de $\mathbf{C}_X$. Si $\mathbf{C}$ est une $S$-categorie abelienne plate, on definit (loc. cit.) une $S$-categorie triangulee $\mathbf{D}(\mathbf{C})$, dont la fibre en $X \in \operatorname{ob} S$ est la categorie derivee $\mathbf{D}(\mathbf{C}_X)$ ; si $f : X \to Y$ est une fleche de $S$, nous noterons $f^* : \mathbf{D}(\mathbf{C}_Y) \to \mathbf{D}(\mathbf{C}_X)$ le foncteur image inverse, qui se deduit de $f^* : \mathbf{K}(\mathbf{C}_Y) \to \mathbf{K}(\mathbf{C}_X)$ par passage au quotient.
02998: \end{definition}
02999: 
03000: \subsubsection*{1.4. Quasi-isomorphismes et conditions de finitude}
03001: \addcontentsline{toc}{subsubsection}{1.4. Quasi-isomorphismes et conditions de finitude}
03002: 
03003: \begin{lemme}[1.4.1]
03004: Soient $S$ un schema, $n \in \mathbb{Z}$, et soit $u : E \to F$ un morphisme de complexes parfaits sur $S$. Supposons que $H^i(u) = 0$ pour $i > n$ et que $H^n(u)$ soit de presentation finie. Alors il existe un diagramme commutatif de complexes parfaits sur $S$
03005: \[
03006: \begin{array}{ccc}
03007: E & \xrightarrow{v} & E' \\
03008: \downarrow{u} & & \downarrow{u'} \\
03009: F & \xrightarrow{=} & F
03010: \end{array}
03011: \]
03012: tel que $H^i(u') = 0$ pour $i \geq n$ et que $v$ soit un isomorphisme en degres $< n$.
03013: \end{lemme}
03014: 
03015: \begin{lemme}[1.4.2]
03016: Soient $\mathbf{C}$ une categorie abelienne, $u : E \to F$ un morphisme de complexes de $\mathbf{C}$ (resp. une fleche de $\mathbf{D}(\mathbf{C})$), $G$ le cone de $u$. Conditions equivalentes pour $n \in \mathbb{Z}$ donne :
03017: \begin{enumerate}
03018: \item[(i)] $H^i(G) = 0$ pour $i \geq n$ ;
03019: \item[(ii)] $H^i(u)$ est un isomorphisme pour $i > n$ et un epimorphisme pour $i = n$.
03020: \end{enumerate}
03021: La preuve est immediate a partir de la suite exacte de cohomologie.
03022: \end{lemme}
03023: 
03024: \begin{definition}[1.4.3]
03025: Nous dirons que $u$ est un \emph{$n$-quasi-isomorphisme} (resp. un \emph{$n$-isomorphisme}) si $u$ verifie les conditions equivalentes de (1.4.2).
03026: \end{definition}
03027: 
03028: Cela etant, on deduit aussitot de (1.4.1) le
03029: 
03030: \begin{corollaire}[1.4.4]
03031: Soient $S$, $\mathbf{C}$, $\mathbf{C}_0$ comme en (1.4.1). Soit $u : E \to F$ un $(n+1)$-quasi-isomorphisme de complexes de $\mathbf{C}_S$ ($n \in \mathbb{Z}$ donne) tel que $H^n(G)$ soit de $\mathbf{C}_0$-type fini. Alors l'ensemble des $X$ au-dessus de $S$ tels qu'il existe $(M, r, s)$ comme en (1.4.1) tels que $u'$ soit un $n$-quasi-isomorphisme est un raffinement de $S$.
03032: \end{corollaire}
03033: 
03034: De maniere imagee, on peut, moyennant la condition de finitude sur $H^n(G)$, localement ``gagner un cran'', i.e. modifier localement $u$ en degre $n$ par l'addition a $E^n$ d'un objet de $\mathbf{C}_0$ de maniere a faire de $u$ un $n$-quasi-isomorphisme.
03035: 
03036: \begin{corollaire}[1.4.5]
03037: Soient $\mathbf{C}$ une categorie abelienne, $n \in \mathbb{Z}$, et $u : E \to F$ un $n$-quasi-isomorphisme de complexes de $\mathbf{C}$. Il existe un diagramme commutatif de $\mathbf{C}(\mathbf{C})$
03038: \[
03039: \begin{array}{ccc}
03040: E & \xrightarrow{v} & E' \\
03041: \downarrow{u} & & \downarrow{u'} \\
03042: F & \xrightarrow{=} & F
03043: \end{array}
03044: \]
03045: tel que $u'$ soit un $n$-isomorphisme, et que $v$ soit un $(n+1)$-quasi-isomorphisme.
03046: \end{corollaire}
03047: 
03048: % [Pages 93-100: More technical results on derived categories]
03049: 
03050: \subsubsection*{1.5. Conditions de finitude dans les categories derivees}
03051: \addcontentsline{toc}{subsubsection}{1.5. Conditions de finitude dans les categories derivees}
03052: 
03053: \begin{definition}
03054: Nous dirons que $\mathbf{C}_0$ est \emph{localement stable par noyau d'epimorphisme} si la condition suivante est realisee : pour tout $X \in \operatorname{ob} S$ et tout epimorphisme $u : E \to F$ de $\mathbf{C}_X$ avec $E$ et $F \in \operatorname{ob} \mathbf{C}_{0,X}$, le noyau de $u$ est localement dans $\mathbf{C}_{0,X}$ (i.e. l'ensemble des $f : U \to X$ de $S$ tels que $f^*(\Ker(u)) \in \operatorname{ob} \mathbf{C}_{0,U}$ est un raffinement de $X$).
03055: \end{definition}
03056: 
03057: Par le procede standard de decomposition d'une suite exacte longue en suites exactes courtes, on voit que la condition precedente equivaut a la condition (en apparence plus forte) que voici : pour tout $X \in \operatorname{ob} S$ et toute suite exacte de $\mathbf{C}_X$
03058: \[
03059: (*) \quad 0 \longrightarrow E^0 \longrightarrow E^1 \longrightarrow \cdots \longrightarrow E^n \longrightarrow 0
03060: \]
03061: avec $E^i \in \operatorname{ob} \mathbf{C}_{0,X}$ pour $1 \leq i \leq n$, $E^0$ est localement dans $\mathbf{C}_{0,X}$.
03062: 
03063: Nous dirons que $\mathbf{C}_0$ est \emph{localement quasi-stable par noyau d'epimorphismes} si, pour toute suite exacte telle que $(*)$, $E^0$ est de $\mathbf{C}_0$-type fini.
03064: 
03065: \begin{definition}
03066: Nous dirons que $\mathbf{C}_0$ est \emph{localement quasi-relevable} dans $\mathbf{C}$ si la condition plus faible suivante est realisee : pour tout $X \in \operatorname{ob} S$, tout epimorphisme $u : E \to F$ de $\mathbf{C}_X$ avec $F \in \operatorname{ob} \mathbf{C}_{0,X}$ et $E \in \operatorname{ob} \mathbf{C}_X$, il existe localement un epimorphisme $v : E' \to E$ avec $E' \in \operatorname{ob} \mathbf{C}_{0,X}$ tel que $uv$ soit un epimorphisme de $\mathbf{C}_{0,X}$.
03067: \end{definition}
03068: 
03069: \begin{theoreme}[1.5.1]
03070: Supposons que $\mathbf{C}_0$ soit localement stable par noyau d'epimorphisme (resp. localement quasi-stable par noyau d'epimorphismes, resp. localement quasi-relevable dans $\mathbf{C}$). Alors la sous-categorie $\mathbf{D}_{\mathbf{C}_0}(\mathbf{C})$ de $\mathbf{D}(\mathbf{C})$ formee des complexes dont la cohomologie est localement dans $\mathbf{C}_0$ (resp. de $\mathbf{C}_0$-type fini, resp. localement quasi-relevable) est une sous-categorie triangulee de $\mathbf{D}(\mathbf{C})$.
03071: \end{theoreme}
03072: 
03073: \vfill
03074: \begin{center}
03075: \rule{0.5\textwidth}{0.4pt}
03076: 
03077: \textit{Fin de la selection des pages 51--100 de SGA 6}
03078: \end{center}
03079: 
03080: 
03081: \clearpage
03082: % ===== Source block: SGA 6 pages 101 150 =====
03083: \phantomsection
03084: \addcontentsline{toc}{section}{SGA 6 pages 101 150}
03085: \section*{SGA 6 pages 101 150}
03086: \setcounter{page}{101}
03087: 
03088: \selectlanguage{french}
03089: 
03090: %====================================================
03091: % PAGE 101 (fin du \S 1)
03092: %====================================================
03093: \section*{Corollaire 1.4.5}